%% ---------------------------------------------------------------------
%% Appendix B -- Mathematical Refresher
%% Every result is stated precisely and tied to the chapters that use it.
%% Cross-references to chapters use ch:chNN only.  Every number quoted in
%% the examples is produced by code/figures/gen_appB_numbers.py.
%% ---------------------------------------------------------------------
\chapter{Mathematical Refresher}
\label{ch:appB}
\index{mathematical refresher|(}

\begin{objectives}
After working through this appendix you will be able to
\begin{itemize}
  \item use the matrix identities (transpose, inverse, eigen-decomposition, definiteness, block matrices, Schur complement, least squares) that the filtering, control and consensus chapters rely on;
  \item linearise a nonlinear map with its Jacobian, bound the error, and step a double integrator forward exactly;
  \item write down a multivariate Gaussian, transform it affinely, condition it on a measurement, and gate with the Mahalanobis distance;
  \item tell a convex problem from a non-convex one, read the KKT conditions, and say what linear, quadratic and mixed-integer programs are and how each is solved;
  \item form the Laplacian of a graph and explain why its second eigenvalue measures connectivity;
  \item read a complexity bound and choose the right \python data structure for a search.
\end{itemize}
\end{objectives}

\noindent
This appendix collects the mathematics that the chapters use without proving
it. It is a refresher, not a course: every result is stated exactly, short
proofs are given in full, long ones are sketched with a pointer to a textbook,
and each result names the chapters that depend on it. Read a section when a
chapter sends you here; nothing has to be read in order.
The notation is that of the notation table at the front of
the book, and every number quoted in the examples is produced by the script
\code{code/figures/gen\_appB\_numbers.py}, so you can check your own hand
calculations against it.

\begin{keyidea}
Four moves carry almost every derivation in this book: diagonalise a symmetric
matrix ($\mat{A} = \mat{Q}\mat{\Lambda}\mat{Q}\T$), linearise a nonlinear map
(keep the Jacobian, bound the rest), push a Gaussian through an affine map and
condition it on what was measured (the Kalman filter is nothing else), and
check convexity before trusting a local minimum.
\end{keyidea}

% ---------------------------------------------------------------------
\section{Vectors and matrices}
\label{sec:appB-linalg}

\subsection{Vectors, norms and products}

A vector $\vect{x} \in \R^n$ is a column of $n$ real numbers $x_1, \dots,
x_n$; its transpose $\vect{x}\T$ is the corresponding row. Positions,
velocities and filter states are vectors of this kind, with $n = 2$ or $3$ for
a point in the plane or in space and a larger $n$ for a stacked state such as
$\state = (\pos, \vel)$.

\begin{definition}[Norms, dot product, cross product]
\label{def:appB-norm}
For $\vect{x}, \vect{y} \in \R^n$ the \textbf{dot product}\index{dot product}
is $\vect{x} \cdot \vect{y} = \vect{x}\T \vect{y} = \sum_i x_i y_i$. The
\textbf{Euclidean norm}\index{norm!Euclidean} is $\norm{\vect{x}} =
\sqrt{\vect{x}\T\vect{x}}$, the \textbf{1-norm} is $\norm{\vect{x}}_1 = \sum_i
\abs{x_i}$ and the \textbf{$\infty$-norm} is $\norm{\vect{x}}_\infty = \max_i
\abs{x_i}$.\index{norm!1-norm and infinity-norm} Every norm satisfies
$\norm{\vect{x}} \ge 0$ with equality only for $\vect{x} = \vect{0}$,
$\norm{\alpha \vect{x}} = \abs{\alpha}\,\norm{\vect{x}}$, and the triangle
inequality $\norm{\vect{x} + \vect{y}} \le \norm{\vect{x}} + \norm{\vect{y}}$.
For $\vect{a}, \vect{b} \in \R^3$ the \textbf{cross product}\index{cross product}
is
\begin{equation}
  \vect{a} \times \vect{b} = \left( a_y b_z - a_z b_y,\; a_z b_x - a_x b_z,\; a_x b_y - a_y b_x \right)\T ,
  \label{eq:appB-cross}
\end{equation}
and for $\vect{a}, \vect{b} \in \R^2$ the scalar $\vect{a} \times \vect{b} =
a_x b_y - a_y b_x = \vect{a}^{\perp} \cdot \vect{b}$, with $\vect{a}^{\perp} =
(-a_y, a_x)\T$ the quarter turn of $\vect{a}$.
\end{definition}

Three facts about these products are used all over the geometry of
\cref{ch:ch02,ch:ch12,ch:ch13}. First, $\vect{x} \cdot \vect{y} = \norm{\vect{x}}
\norm{\vect{y}} \cos\theta$ with $\theta$ the angle between the vectors, so
$\abs{\vect{x} \cdot \vect{y}} \le \norm{\vect{x}} \norm{\vect{y}}$ (the
\textbf{Cauchy--Schwarz inequality}\index{Cauchy--Schwarz inequality}) and two
vectors are orthogonal exactly when their dot product is zero. Second, the
\textbf{projection} of $\vect{x}$ onto a unit vector $\vect{u}$ is $(\vect{x}
\cdot \vect{u})\, \vect{u}$, which is how the distance from a point to a
segment is computed in \cref{ch:ch02}. Third, $\vect{a} \times \vect{b}$ in
$\R^3$ is orthogonal to both factors with length $\norm{\vect{a}}
\norm{\vect{b}} \sin\theta$, and the planar scalar $\vect{a} \times \vect{b}$
is positive exactly when $\vect{b}$ lies counter-clockwise from $\vect{a}$:
this sign is the left-or-right test behind the half-plane and cone tests of
\cref{ch:ch12,ch:ch13}. The grid heuristics of \cref{ch:ch04} are norms of the
displacement $\vect{d}$ to the goal: Manhattan $\norm{\vect{d}}_1$, Euclidean
$\norm{\vect{d}}$, Chebyshev $\norm{\vect{d}}_\infty$.

\subsection{Matrices, transpose, inverse and rank}

A matrix $\mat{A} \in \R^{m \times n}$ maps $\R^n$ to $\R^m$ by $\vect{x}
\mapsto \mat{A}\vect{x}$; the product $\mat{A}\mat{B}$, with entries
$(\mat{A}\mat{B})_{ij} = \sum_k a_{ik} b_{kj}$, is the composition ``first
$\mat{B}$, then $\mat{A}$'' and is not commutative in general. The
\textbf{transpose}\index{matrix!transpose} $\mat{A}\T$ swaps rows and columns
(a square $\mat{A}$ with $\mat{A}\T = \mat{A}$ is \textbf{symmetric}); the
\textbf{inverse}\index{matrix!inverse} $\mat{A}\inv$ of a square matrix
satisfies $\mat{A}\mat{A}\inv = \mat{A}\inv\mat{A} = \mat{I}$, when it exists;
the \textbf{rank}\index{matrix!rank} is the number of linearly independent
columns, which equals the number of linearly independent rows; the
\textbf{trace} $\tr \mat{A} = \sum_i a_{ii}$ and the \textbf{determinant}
$\det \mat{A}$ are defined for square matrices; and an \textbf{orthogonal}
matrix, $\mat{Q}\T\mat{Q} = \mat{I}$, such as a rotation, preserves norms and
dot products. The following equivalences are the working definition of
invertibility.

\begin{proposition}[Invertibility]
\label{thm:appB-invertible}
For a square matrix $\mat{A} \in \R^{n \times n}$ the following are
equivalent: $\mat{A}$ is invertible; $\operatorname{rank} \mat{A} = n$;
$\det \mat{A} \ne 0$; $\mat{A}\vect{x} = \vect{0}$ only for $\vect{x} =
\vect{0}$; $\mat{A}$ has no zero eigenvalue. In that case
$\mat{A}\vect{x} = \vect{b}$ has the unique solution $\vect{x} =
\mat{A}\inv\vect{b}$.
\end{proposition}

\Cref{tab:appB-identities} lists the algebra of transposes, inverses, traces
and determinants that the derivations of \cref{ch:ch18,ch:ch21,ch:ch23} use
without comment; the block-matrix rows are explained in
\cref{sec:appB-schur}.

\begin{table}[tb]
  \centering\small
  \caption{Matrix identities used in the book. All inverses are assumed to exist; $\lambda_i$ are the eigenvalues; $\mat{S} = \mat{A} - \mat{B}\mat{D}\inv\mat{C}$ is the Schur complement of $\mat{D}$ in the block matrix $\mat{M}$ of \cref{thm:appB-block}.}
  \label{tab:appB-identities}
  \begin{tabular}{@{}L{8.9cm}L{5.9cm}@{}}
  \toprule
  Identity & Remark\\
  \midrule
  $(\mat{A}\mat{B})\T = \mat{B}\T\mat{A}\T$, \ $(\mat{A}\mat{B})\inv = \mat{B}\inv\mat{A}\inv$ & the order reverses\\
  $(\mat{A}\T)\inv = (\mat{A}\inv)\T$ & written $\mat{A}^{-\mathsf{T}}$\\
  $\tr(\mat{A}\mat{B}) = \tr(\mat{B}\mat{A})$, \ $\tr \mat{A} = \sum_i \lambda_i$ & the trace is cyclic\\
  $\det(\mat{A}\mat{B}) = \det\mat{A}\,\det\mat{B}$, \ $\det\mat{A}\T = \det\mat{A}$, \ $\det\mat{A} = \prod_i \lambda_i$ & $\det \mat{A}\inv = 1/\det \mat{A}$\\
  $\vect{x}\T\mat{A}\vect{x} = \tfrac12\vect{x}\T(\mat{A} + \mat{A}\T)\vect{x}$ & only the symmetric part matters\\
  $\mat{A} = \mat{Q}\mat{\Lambda}\mat{Q}\T$, \ $\mat{A}\inv = \mat{Q}\mat{\Lambda}\inv\mat{Q}\T$ \ ($\mat{A}$ symmetric) & \cref{thm:appB-spectral}\\
  $\mat{A} = \mat{L}\mat{L}\T$ \ ($\mat{A} \succ \mat{0}$, $\mat{L}$ lower triangular) & Cholesky factorisation\\
  $\det\mat{M} = \det\mat{S}\,\det\mat{D}$, \ $(\mat{M}\inv)_{11} = \mat{S}\inv$ & block matrices, \cref{thm:appB-block}\\
  $(\mat{A} + \mat{U}\mat{C}\mat{V})\inv = \mat{A}\inv - \mat{A}\inv\mat{U}(\mat{C}\inv + \mat{V}\mat{A}\inv\mat{U})\inv\mat{V}\mat{A}\inv$ & Woodbury identity, \cref{thm:appB-woodbury}\\
  $\nabla_{\vect{x}}(\vect{b}\T\vect{x}) = \vect{b}$, \ $\nabla_{\vect{x}}(\vect{x}\T\mat{A}\vect{x}) = (\mat{A} + \mat{A}\T)\vect{x}$ & derivatives, \cref{sec:appB-calculus}\\
  \bottomrule
  \end{tabular}
\end{table}

\subsection{Eigenvalues, symmetric matrices and definiteness}
\label{sec:appB-psd}

\begin{definition}[Eigenvalues and eigenvectors]
\label{def:appB-eigen}
A scalar $\lambda$ and a non-zero vector $\vect{q}$ with $\mat{A}\vect{q} =
\lambda\vect{q}$ are an \textbf{eigenvalue}\index{eigenvalue} and an
\textbf{eigenvector} of the square matrix $\mat{A}$. The eigenvalues are the
roots of $\det(\mat{A} - \lambda\mat{I}) = 0$.
\end{definition}

For a general real matrix the eigenvalues may be complex. The matrices that
carry uncertainty, cost and connectivity in this book are symmetric, and for
them everything is real and orthogonal.

\begin{theorem}[Spectral theorem]
\label{thm:appB-spectral}
A symmetric matrix $\mat{A} \in \R^{n \times n}$ has $n$ real eigenvalues
$\lambda_1 \ge \dots \ge \lambda_n$ with orthonormal eigenvectors
$\vect{q}_1, \dots, \vect{q}_n$, so that
\begin{equation}
  \mat{A} = \mat{Q}\mat{\Lambda}\mat{Q}\T = \sum_{i=1}^{n} \lambda_i\, \vect{q}_i \vect{q}_i\T,
  \qquad \mat{Q} = [\vect{q}_1 \cdots \vect{q}_n],\ \mat{Q}\T\mat{Q} = \mat{I},\
  \mat{\Lambda} = \diag(\lambda_1, \dots, \lambda_n).
  \label{eq:appB-spectral}
\end{equation}
\end{theorem}

The proof is in every linear-algebra text \cite[Appendix~A]{boyd2004convex}.
The decomposition says that a symmetric matrix is a pure scaling by
$\lambda_i$ along each of $n$ orthogonal axes $\vect{q}_i$; hence $\mat{A}^k =
\mat{Q}\mat{\Lambda}^k\mat{Q}\T$ and, if no eigenvalue vanishes, $\mat{A}\inv =
\mat{Q}\mat{\Lambda}\inv\mat{Q}\T$. The \textbf{quadratic form}\index{quadratic form}
$\vect{x}\T\mat{A}\vect{x} = \sum_i \lambda_i (\vect{q}_i\T\vect{x})^2$ is a
weighted sum of squared coordinates along the axes, and its sign decides
whether a covariance is valid, a cost is convex and a consensus protocol
converges.

\begin{definition}[Positive definite and semidefinite matrices]
\label{def:appB-psd}
A symmetric matrix $\mat{A}$ is \textbf{positive semidefinite}\index{matrix!positive semidefinite}
(PSD), written $\mat{A} \succeq \mat{0}$, if $\vect{x}\T\mat{A}\vect{x} \ge 0$
for all $\vect{x}$, and \textbf{positive definite}\index{matrix!positive definite}
(PD), written $\mat{A} \succ \mat{0}$, if $\vect{x}\T\mat{A}\vect{x} > 0$ for
all $\vect{x} \ne \vect{0}$. $\mat{A} \succeq \mat{B}$ means $\mat{A} - \mat{B}
\succeq \mat{0}$.
\end{definition}

\begin{proposition}[Tests and closure properties]
\label{thm:appB-psd-tests}
Let $\mat{A}$ be symmetric.
\begin{enumerate}[label=(\roman*)]
  \item $\mat{A} \succeq \mat{0}$ iff all eigenvalues are $\ge 0$, and $\mat{A} \succ \mat{0}$ iff all are $> 0$. A PD matrix is invertible and its inverse is PD.
  \item $\mat{A} \succeq \mat{0}$ iff $\mat{A} = \mat{B}\T\mat{B}$ for some $\mat{B}$. If $\mat{A} \succ \mat{0}$, $\mat{B}\T$ can be chosen lower triangular with positive diagonal, $\mat{A} = \mat{L}\mat{L}\T$ (the \textbf{Cholesky factorisation}\index{Cholesky factorisation}), and this $\mat{L}$ is unique.
  \item If $\mat{A}, \mat{B} \succeq \mat{0}$ and $\alpha, \beta \ge 0$ then $\alpha\mat{A} + \beta\mat{B} \succeq \mat{0}$; for any $\mat{C}$ of compatible size $\mat{C}\T\mat{A}\mat{C} \succeq \mat{0}$; and $\mat{A} + \eps\mat{I} \succ \mat{0}$ for every $\eps > 0$.
\end{enumerate}
\end{proposition}
\begin{proof}[Proof sketch]
(i) and (iii) are read off the quadratic form $\sum_i \lambda_i (\vect{q}_i\T\vect{x})^2$ and the definition, using $\vect{x}\T\mat{C}\T\mat{A}\mat{C}\vect{x} = (\mat{C}\vect{x})\T\mat{A}(\mat{C}\vect{x})$; the eigenvalues of $\mat{A}\inv$ are $1/\lambda_i$. For (ii), $\vect{x}\T\mat{B}\T\mat{B}\vect{x} = \norm{\mat{B}\vect{x}}^2 \ge 0$; conversely $\mat{B} = \mat{\Lambda}^{1/2}\mat{Q}\T$ works, and Gaussian elimination without pivoting produces the unique triangular factor \cite[Appendix~C]{boyd2004convex}. \qedhere
\end{proof}

Geometrically, the level set $\vect{x}\T\mat{A}\vect{x} = 1$ of a PD quadratic
form is an ellipsoid whose axes are the eigenvectors $\vect{q}_i$ and whose
semi-axes are $1/\sqrt{\lambda_i}$; the covariance ellipse of a Gaussian in
\cref{sec:appB-probability} is the level set of the form built from
$\mat{\Sigma}\inv$, so its semi-axes are $\sqrt{\lambda_i(\mat{\Sigma})}$.

\begin{example}[A covariance matrix taken apart]
\label{ex:appB-eigen}
The matrix
\begin{equation}
  \mat{\Sigma} = \begin{bmatrix} 2.0 & 1.2\\ 1.2 & 1.0 \end{bmatrix}
  \label{eq:appB-sigma}
\end{equation}
is symmetric with $\tr\mat{\Sigma} = 3$ and $\det\mat{\Sigma} = 0.56$. Its
eigenvalues solve $\lambda^2 - 3\lambda + 0.56 = 0$ and are $\lambda_1 = 2.8$
and $\lambda_2 = 0.2$: both positive, so $\mat{\Sigma} \succ \mat{0}$. The
eigenvector of $\lambda_1$ is $\vect{q}_1 = (3, 2)\T/\sqrt{13} = (0.832,
0.555)\T$, at an angle of $33.7^\circ$ to the $x$-axis, and $\vect{q}_2 =
(-2, 3)\T/\sqrt{13}$ is orthogonal to it. The ellipse
$\vect{x}\T\mat{\Sigma}\inv\vect{x} = 1$ has semi-axes $\sqrt{2.8} = 1.673$
along $\vect{q}_1$ and $\sqrt{0.2} = 0.447$ along $\vect{q}_2$; it is drawn in
\cref{fig:appB-conditioning}. The Cholesky factor is
$\mat{L} = \begin{bsmallmatrix} 1.4142 & 0\\ 0.8485 & 0.5292\end{bsmallmatrix}$,
and $\mat{\Sigma}\inv = \frac{1}{0.56}\begin{bsmallmatrix} 1.0 & -1.2\\ -1.2 &
2.0\end{bsmallmatrix} = \begin{bsmallmatrix} 1.786 & -2.143\\ -2.143 &
3.571\end{bsmallmatrix}$. \Cref{sec:appB-probability} uses this
$\mat{\Sigma}$ as the covariance of a Gaussian.
\end{example}

\subsection{Block matrices and the Schur complement}
\label{sec:appB-schur}

The large matrices of this book come in blocks: the covariance of $(\pos,
\vel)$, the joint covariance of a state and a measurement, the KKT matrix of a
quadratic program. Blocks multiply like scalars as long as their order is kept,
\begin{equation}
  \begin{bmatrix}\mat{A} & \mat{B}\\ \mat{C} & \mat{D}\end{bmatrix}
  \begin{bmatrix}\mat{E} & \mat{F}\\ \mat{G} & \mat{H}\end{bmatrix} =
  \begin{bmatrix}\mat{A}\mat{E} + \mat{B}\mat{G} & \mat{A}\mat{F} + \mat{B}\mat{H}\\
                 \mat{C}\mat{E} + \mat{D}\mat{G} & \mat{C}\mat{F} + \mat{D}\mat{H}\end{bmatrix},
  \label{eq:appB-block-product}
\end{equation}
a block-diagonal matrix is inverted block by block, and the transpose of a
block matrix transposes both the arrangement and the blocks.

\begin{definition}[Schur complement]
\label{def:appB-schur}
Let $\mat{M} = \begin{bsmallmatrix}\mat{A} & \mat{B}\\ \mat{C} & \mat{D}\end{bsmallmatrix}$ with $\mat{A}$ and $\mat{D}$ square. If $\mat{D}$ is invertible, the \textbf{Schur complement}\index{Schur complement} of $\mat{D}$ in $\mat{M}$ is $\mat{S} = \mat{A} - \mat{B}\mat{D}\inv\mat{C}$; if $\mat{A}$ is invertible, the Schur complement of $\mat{A}$ in $\mat{M}$ is $\mat{S}' = \mat{D} - \mat{C}\mat{A}\inv\mat{B}$.
\end{definition}

\begin{theorem}[Block factorisation, determinant and inverse]
\label{thm:appB-block}
If $\mat{D}$ is invertible and $\mat{S} = \mat{A} - \mat{B}\mat{D}\inv\mat{C}$, then
\begin{equation}
  \mat{M} =
  \begin{bmatrix}\mat{I} & \mat{B}\mat{D}\inv\\ \mat{0} & \mat{I}\end{bmatrix}
  \begin{bmatrix}\mat{S} & \mat{0}\\ \mat{0} & \mat{D}\end{bmatrix}
  \begin{bmatrix}\mat{I} & \mat{0}\\ \mat{D}\inv\mat{C} & \mat{I}\end{bmatrix},
  \qquad \det\mat{M} = \det\mat{S}\,\det\mat{D},
  \label{eq:appB-block-fact}
\end{equation}
and $\mat{M}$ is invertible iff $\mat{S}$ is, with
\begin{equation}
  \mat{M}\inv =
  \begin{bmatrix}\mat{S}\inv & -\mat{S}\inv\mat{B}\mat{D}\inv\\
   -\mat{D}\inv\mat{C}\mat{S}\inv & \mat{D}\inv + \mat{D}\inv\mat{C}\mat{S}\inv\mat{B}\mat{D}\inv\end{bmatrix}.
  \label{eq:appB-block-inverse}
\end{equation}
If instead $\mat{A}$ is invertible and $\mat{S}' = \mat{D} - \mat{C}\mat{A}\inv\mat{B}$, the same statements hold with the roles of the blocks exchanged; in particular $\det\mat{M} = \det\mat{A}\,\det\mat{S}'$, $(\mat{M}\inv)_{22} = \mat{S}'^{-1}$ and $(\mat{M}\inv)_{11} = \mat{A}\inv + \mat{A}\inv\mat{B}\mat{S}'^{-1}\mat{C}\mat{A}\inv$.
\end{theorem}
\begin{proof}
Multiplying out the right-hand side of \cref{eq:appB-block-fact} with \cref{eq:appB-block-product} gives back $\mat{M}$. The outer factors are block-triangular with identity diagonal blocks, so they have determinant $1$ and are inverted by negating the off-diagonal block; the middle factor has determinant $\det\mat{S}\det\mat{D}$ and is inverted block by block. Inverting the three factors in reverse order and multiplying out gives \cref{eq:appB-block-inverse}. The second statement is the first applied to $\mat{M}$ with its block rows and block columns swapped. \qedhere
\end{proof}

In \cref{ex:appB-eigen} the Schur complement of the $(2,2)$ entry is $2.0 -
1.2^2/1.0 = 0.56$, which is both $\det\mat{\Sigma}/\Sigma_{22}$ and the
reciprocal of the top-left entry $1.786$ of $\mat{\Sigma}\inv$, as the theorem
says. Reading the theorem in both directions gives the identity with which
\cref{ch:ch18} passes between the covariance and the information form of the
Kalman update, and which turns the inversion of a large matrix into that of a
small one whenever the update has low rank.

\begin{proposition}[Matrix inversion lemma, Woodbury identity]
\label{thm:appB-woodbury}
For $\mat{A} \in \R^{n\times n}$, $\mat{U} \in \R^{n \times k}$, $\mat{C} \in \R^{k \times k}$ and $\mat{V} \in \R^{k \times n}$, whenever all inverses below exist,
\begin{equation}
  (\mat{A} + \mat{U}\mat{C}\mat{V})\inv = \mat{A}\inv - \mat{A}\inv\mat{U}\left(\mat{C}\inv + \mat{V}\mat{A}\inv\mat{U}\right)\inv\mat{V}\mat{A}\inv .
  \label{eq:appB-woodbury}
\end{equation}
\end{proposition}
\begin{proof}
Apply \cref{thm:appB-block} to $\mat{M} = \begin{bsmallmatrix}\mat{A} & \mat{U}\\ \mat{V} & -\mat{C}\inv\end{bsmallmatrix}$: eliminating the $(2,2)$ block gives $(\mat{M}\inv)_{11} = (\mat{A} + \mat{U}\mat{C}\mat{V})\inv$, eliminating the $(1,1)$ block gives $(\mat{M}\inv)_{11} = \mat{A}\inv - \mat{A}\inv\mat{U}(\mat{C}\inv + \mat{V}\mat{A}\inv\mat{U})\inv\mat{V}\mat{A}\inv$, and the two expressions are equal. \qedhere
\end{proof}

\subsection{Least squares and the normal equations}
\label{sec:appB-lsq}

\begin{theorem}[Normal equations]
\label{thm:appB-normal}
Let $\mat{A} \in \R^{m \times n}$ and $\vect{b} \in \R^m$. A vector $\vect{x}^*$ minimises $\norm{\mat{A}\vect{x} - \vect{b}}^2$ iff it satisfies the \textbf{normal equations}\index{least squares}\index{normal equations}
\begin{equation}
  \mat{A}\T\mat{A}\,\vect{x}^* = \mat{A}\T\vect{b}.
  \label{eq:appB-normal}
\end{equation}
If $\mat{A}$ has full column rank $n$, then $\mat{A}\T\mat{A} \succ \mat{0}$ and the minimiser is unique, $\vect{x}^* = (\mat{A}\T\mat{A})\inv\mat{A}\T\vect{b}$. For a weight matrix $\mat{W} \succ \mat{0}$ the minimiser of $(\mat{A}\vect{x} - \vect{b})\T\mat{W}(\mat{A}\vect{x} - \vect{b})$ solves $\mat{A}\T\mat{W}\mat{A}\,\vect{x}^* = \mat{A}\T\mat{W}\vect{b}$.
\end{theorem}
\begin{proof}
With $\vect{r} = \mat{A}\vect{x}^* - \vect{b}$ and $\vect{d} = \vect{x} - \vect{x}^*$, $\norm{\mat{A}\vect{x} - \vect{b}}^2 = \norm{\vect{r}}^2 + 2\vect{d}\T\mat{A}\T\vect{r} + \norm{\mat{A}\vect{d}}^2$. If \cref{eq:appB-normal} holds, the middle term vanishes for every $\vect{d}$ and $\vect{x}^*$ is a minimiser; if not, $\vect{d} = -t\mat{A}\T\vect{r}$ with a small $t > 0$ makes the last two terms sum to a negative number. Full column rank gives $\norm{\mat{A}\vect{d}}^2 > 0$ for $\vect{d} \ne \vect{0}$, hence uniqueness; the weighted case is the unweighted one for $\mat{W}^{1/2}\mat{A}$ and $\mat{W}^{1/2}\vect{b}$, with $\mat{W}^{1/2} = \mat{Q}\mat{\Lambda}^{1/2}\mat{Q}\T$ from \cref{thm:appB-spectral}. \qedhere
\end{proof}

\begin{example}[Fitting a constant-velocity line]
\label{ex:appB-lsq}
A drone is seen at the positions $0.0$, $1.1$ and $1.9$\,m at the times $t = 0, 1, 2$\,s along one axis. The model $p(t) = p_0 + vt$ is linear in $\vect{x} = (p_0, v)\T$, with rows $(1, t)$ in $\mat{A}$ and $\vect{b} = (0.0, 1.1, 1.9)\T$, so $\mat{A}\T\mat{A} = \begin{bsmallmatrix} 3 & 3\\ 3 & 5\end{bsmallmatrix}$, $\mat{A}\T\vect{b} = (3.0, 4.9)\T$, and \cref{eq:appB-normal} gives $p_0 = 0.05$\,m and $v = 0.95$\,m/s. The residuals $\mat{A}\vect{x}^* - \vect{b} = (0.05, -0.10, 0.05)\T$ have squared norm $0.015$ and are orthogonal to both columns of $\mat{A}$, which is what the normal equations say.
\end{example}

Least squares is the quiet workhorse of Parts VI and VII: the Kalman update
of \cref{ch:ch18} is a weighted least-squares compromise between a prediction
and a measurement, weighted by the inverses of their covariances, and an MPC
problem without inequality constraints (\cref{ch:ch21}) is a least-squares
problem in the stacked inputs, which the constraints turn into the quadratic
program of \cref{sec:appB-qp}.

\begin{pitfall}[Do not form the inverse]
The formulas above are written with $\mat{A}\inv$ because that is how they are proved, not how they should be computed. Solving $\mat{A}\vect{x} = \vect{b}$ with \code{numpy.linalg.solve}, or with a Cholesky factor when $\mat{A} \succ \mat{0}$, is faster and more accurate than forming \code{numpy.linalg.inv(A)} and multiplying. Two more habits pay off in \cref{ch:ch18,ch:ch19}: re-symmetrise a covariance as $\tfrac12(\mat{P} + \mat{P}\T)$ after each update, and check its smallest eigenvalue before a Cholesky factorisation, because rounding can turn a PSD matrix into one with a slightly negative eigenvalue.
\end{pitfall}

% ---------------------------------------------------------------------
\section{Calculus for motion}
\label{sec:appB-calculus}

\subsection{Gradient, Jacobian, Hessian and the chain rule}

\begin{definition}[Derivatives of vector functions]
\label{def:appB-gradient}
Let $f\colon \R^n \to \R$ and $\vect{g}\colon \R^n \to \R^m$ be differentiable. The \textbf{gradient}\index{gradient} $\nabla f(\vect{x}) \in \R^n$ has the entries $\partial f/\partial x_i$; the \textbf{Jacobian}\index{Jacobian} $\mat{J}_{\vect{g}}(\vect{x}) \in \R^{m \times n}$ has the entries $\partial g_i/\partial x_j$, one row per output and one column per input, so that the gradient is the transpose of the Jacobian of a scalar function; the \textbf{Hessian}\index{Hessian} $\nabla^2 f(\vect{x}) \in \R^{n \times n}$ has the entries $\partial^2 f/\partial x_i \partial x_j$ and is symmetric when $f$ is twice continuously differentiable.
\end{definition}

The gradient points in the direction of steepest ascent and is orthogonal to
the level set of $f$ through $\vect{x}$; the derivative of $f$ along a unit
vector $\vect{u}$ is $\nabla f(\vect{x})\T\vect{u}$. The \textbf{chain
rule}\index{chain rule} for a composition is a matrix product: if
$\vect{h}(\vect{x}) = \vect{g}(\vect{f}(\vect{x}))$ then
\begin{equation}
  \mat{J}_{\vect{h}}(\vect{x}) = \mat{J}_{\vect{g}}(\vect{f}(\vect{x}))\,\mat{J}_{\vect{f}}(\vect{x}),
  \label{eq:appB-chain}
\end{equation}
which is how the gradients of a neural network in \cref{ch:ch20} are
propagated backwards through its layers. The derivatives that recur in this
book are
\begin{equation}
  \begin{aligned}
  \nabla(\vect{b}\T\vect{x}) &= \vect{b}, &
  \nabla\left(\tfrac12\vect{x}\T\mat{A}\vect{x}\right) &= \mat{A}\vect{x}\ \ (\mat{A} \text{ symmetric}),\\
  \nabla\left(\tfrac12\norm{\vect{x} - \vect{c}}^2\right) &= \vect{x} - \vect{c}, &
  \nabla\norm{\vect{x} - \vect{c}} &= \frac{\vect{x} - \vect{c}}{\norm{\vect{x} - \vect{c}}},
  \end{aligned}
  \label{eq:appB-derivs}
\end{equation}
together with $\mat{J}_{\mat{A}\vect{x}} = \mat{A}$ and
$\nabla^2(\tfrac12\vect{x}\T\mat{A}\vect{x}) = \mat{A}$. The attractive
potential of \cref{ch:ch15} is the third expression with $\vect{c}$ the goal,
so its force $-\nabla U$ points at the goal; the repulsive potential is a
function of the clearance to an obstacle, and by the chain rule its gradient
is a multiple of the unit vector in the fourth expression, which points away
from the obstacle.

\begin{proposition}[Taylor expansion]
\label{thm:appB-taylor}
If $\vect{g}\colon \R^n \to \R^m$ is twice continuously differentiable near $\vect{x}$, then for small $\vect{\delta}$
\begin{equation}
  \vect{g}(\vect{x} + \vect{\delta}) = \vect{g}(\vect{x}) + \mat{J}_{\vect{g}}(\vect{x})\,\vect{\delta} + \vect{r}(\vect{\delta}),\qquad \norm{\vect{r}(\vect{\delta})} \le c\,\norm{\vect{\delta}}^2
  \label{eq:appB-taylor1}
\end{equation}
for a constant $c$ that depends on the second derivatives of $\vect{g}$ near $\vect{x}$. For a scalar $f$ the next term is explicit: $f(\vect{x} + \vect{\delta}) = f(\vect{x}) + \nabla f(\vect{x})\T\vect{\delta} + \tfrac12\vect{\delta}\T\nabla^2 f(\vect{x})\vect{\delta} + \bigO{\norm{\vect{\delta}}^3}$.
\end{proposition}

Dropping $\vect{r}$ in \cref{eq:appB-taylor1} is
\textbf{linearisation}\index{linearisation}. The extended Kalman filter of
\cref{ch:ch19} replaces its motion and measurement models by their Jacobians
$\mat{F}$ and $\mat{H}$ at the current estimate, and the size of $\vect{r}$
over the spread of the estimate decides whether that is good enough or whether
the unscented or particle filters of the same chapter are needed. The
second-order expansion underlies Newton's method and the convexity test of
\cref{sec:appB-optimization}.

\begin{example}[Range and bearing of a target]
\label{ex:appB-jacobian}
A radar at the origin measures the range $r$ and the bearing $\varphi$ of a drone with state $\state = (p_x, p_y, v_x, v_y)\T$:
\begin{equation*}
  \vect{h}(\state) = \begin{bmatrix} \sqrt{p_x^2 + p_y^2}\\ \operatorname{atan2}(p_y, p_x)\end{bmatrix},\qquad
  \mat{H} = \mat{J}_{\vect{h}}(\state) = \begin{bmatrix} p_x/r & p_y/r & 0 & 0\\ -p_y/r^2 & p_x/r^2 & 0 & 0\end{bmatrix}.
\end{equation*}
At $p_x = 3$, $p_y = 4$ (so $r = 5$) the Jacobian is $\mat{H} = \begin{bsmallmatrix} 0.6 & 0.8 & 0 & 0\\ -0.16 & 0.12 & 0 & 0\end{bsmallmatrix}$ and $\vect{h} = (5, 0.9273)\T$. For the displacement $\vect{\delta} = (0.1, -0.2, 0, 0)\T$ the linear prediction $\vect{h} + \mat{H}\vect{\delta} = (4.9000, 0.8873)\T$ differs from the exact value $(4.9041, 0.8865)\T$ by $(0.0041, -0.0008)\T$: an error of second order in $\norm{\vect{\delta}} = 0.22$, as \cref{eq:appB-taylor1} promises. The zero columns say that the measurement carries no direct information about the velocity; a filter learns the velocity only through the motion model.
\end{example}

\begin{pitfall}[Check every Jacobian numerically]
Hand-derived Jacobians are the most common source of silent errors in an extended Kalman filter. Compare each column against the central difference $(\vect{h}(\state + \eps\vect{e}_j) - \vect{h}(\state - \eps\vect{e}_j))/2\eps$ with $\eps \approx 10^{-6}$ before you trust it; the script of this appendix does exactly that for \cref{ex:appB-jacobian}. Two details bite in practice: \code{atan2} takes the $y$ argument first, and a bearing difference must be wrapped into $(-\pi, \pi]$ before it is used as an innovation, because a linearisation across the $\pm\pi$ cut is meaningless.
\end{pitfall}

\subsection{The double-integrator step}
\label{sec:appB-double-integrator}

The drone model of \cref{ch:ch02} treats the acceleration $\acc$ as the input
and integrates it twice. If $\acc$ is held constant during a step of length
$\dt$, the velocity grows linearly and the position quadratically,
\begin{equation}
  \vel(t + \dt) = \vel(t) + \acc\,\dt,\qquad
  \pos(t + \dt) = \pos(t) + \vel(t)\,\dt + \tfrac12\acc\,\dt^2 ,
  \label{eq:appB-double-integrator}
\end{equation}
because $\pos(t + \dt) = \pos(t) + \int_0^{\dt}(\vel(t) + \acc\,s)\,\mathrm{d}s$.
This is exact, not an approximation, for a piecewise-constant input, which is
what a digital controller produces. With the state $\state = (\pos, \vel) \in
\R^{2d}$ and the input $\ctrl = \acc \in \R^d$ the step is the linear system
$\state_{k+1} = \mat{A}\state_k + \mat{B}\ctrl_k$ with
\begin{equation}
  \mat{A} = \begin{bmatrix}\mat{I} & \dt\,\mat{I}\\ \mat{0} & \mat{I}\end{bmatrix},\qquad
  \mat{B} = \begin{bmatrix}\tfrac12\dt^2\,\mat{I}\\ \dt\,\mat{I}\end{bmatrix},
  \label{eq:appB-AB}
\end{equation}
where $\mat{I}$ is the $d \times d$ identity.\index{double integrator}
\Cref{ch:ch18} takes $\mat{A}$ as the transition matrix $\mat{F}$ of a
constant-velocity target and turns the unknown acceleration into process noise
through $\mat{B}$, which is where its white-noise-acceleration covariance
$\mat{Q}$ comes from; \cref{ch:ch21} stacks $N$ copies of \cref{eq:appB-AB}
into the prediction matrices of its quadratic program; \cref{ch:ch14} bounds
the velocities reachable in one step by $\norm{\vel_{k+1} - \vel_k} \le
a_{\max}\dt$, its dynamic window. The explicit Euler step $\pos_{k+1} = \pos_k
+ \vel_k\dt$ drops the term $\tfrac12\acc\dt^2$: with $\dt = 0.1$\,s, $v =
2$\,m/s and $a = 1$\,m/s$^2$ the exact step moves the drone $0.205$\,m and
leaves it at $2.1$\,m/s, the Euler step moves it $0.200$\,m, and the difference
accumulates over a trajectory.

% ---------------------------------------------------------------------
\section{Probability and Gaussians}
\label{sec:appB-probability}

\subsection{Random vectors, expectation and covariance}

A \textbf{random vector}\index{random vector} $\state \in \R^n$ is a vector
whose entries are random variables; it is described by a density $p(\state)$
with $p \ge 0$ and $\int p(\state)\,\mathrm{d}\state = 1$, and the probability
that $\state$ falls in a region $\mathcal{R}$ is $\Prob(\state \in \mathcal{R})
= \int_{\mathcal{R}} p(\state)\,\mathrm{d}\state$.

\begin{definition}[Expectation, covariance, cross-covariance]
\label{def:appB-covariance}
The \textbf{expectation}\index{expectation} of $\state$ is $\vect{\mu} = \E[\state] = \int \state\,p(\state)\,\mathrm{d}\state$, taken entry by entry, and $\E[\vect{g}(\state)] = \int \vect{g}(\state)p(\state)\,\mathrm{d}\state$ for a function $\vect{g}$. The \textbf{covariance}\index{covariance} is
\begin{equation}
  \Cov(\state) = \E\left[(\state - \vect{\mu})(\state - \vect{\mu})\T\right] \in \R^{n \times n},
  \label{eq:appB-cov}
\end{equation}
whose diagonal entries are the variances $\Var(x_i) = \sigma_i^2$ and whose off-diagonal entries $\sigma_{ij}$ say how $x_i$ and $x_j$ vary together; $\rho_{ij} = \sigma_{ij}/(\sigma_i\sigma_j) \in [-1, 1]$ is the correlation coefficient. The \textbf{cross-covariance} of two random vectors is $\Cov(\state, \vect{y}) = \E[(\state - \vect{\mu}_x)(\vect{y} - \vect{\mu}_y)\T]$, and $\state$, $\vect{y}$ are \textbf{independent} if $p(\state, \vect{y}) = p(\state)p(\vect{y})$, in which case $\Cov(\state, \vect{y}) = \mat{0}$.
\end{definition}

\begin{proposition}[Rules for means and covariances]
\label{thm:appB-cov-rules}
Expectation is linear, $\E[\mat{A}\state + \mat{B}\vect{y} + \vect{c}] = \mat{A}\E[\state] + \mat{B}\E[\vect{y}] + \vect{c}$, with no assumption about dependence. Every covariance matrix is symmetric and PSD. For fixed $\mat{A}$ and $\vect{b}$,
\begin{equation}
  \Cov(\mat{A}\state + \vect{b}) = \mat{A}\Cov(\state)\mat{A}\T,
  \label{eq:appB-cov-affine}
\end{equation}
and if $\state$ and $\vect{y}$ are independent, $\Cov(\state + \vect{y}) = \Cov(\state) + \Cov(\vect{y})$.
\end{proposition}
\begin{proof}
Linearity is linearity of the integral. For any fixed $\vect{a}$, $\vect{a}\T\Cov(\state)\vect{a} = \E[(\vect{a}\T(\state - \vect{\mu}))^2] = \Var(\vect{a}\T\state) \ge 0$, which is the definition of PSD; symmetry is read off \cref{eq:appB-cov}. Since $\mat{A}\state + \vect{b} - \E[\mat{A}\state + \vect{b}] = \mat{A}(\state - \vect{\mu})$, \cref{eq:appB-cov-affine} follows by pulling $\mat{A}$ and $\mat{A}\T$ out of the expectation. For the sum, expand $(\state + \vect{y} - \vect{\mu}_x - \vect{\mu}_y)(\state + \vect{y} - \vect{\mu}_x - \vect{\mu}_y)\T$; the two cross terms are the cross-covariances, which vanish under independence. \qedhere
\end{proof}

\Cref{eq:appB-cov-affine} is the single most used formula of
\cref{ch:ch18,ch:ch19}: pushing a covariance through the linear model
$\mat{F}$ gives $\mat{F}\mat{P}\mat{F}\T$, and adding independent process
noise adds $\mat{Q}$. Note that it needs no Gaussian assumption at all.

\subsection{The multivariate Gaussian}

\begin{definition}[Multivariate Gaussian]
\label{def:appB-gaussian}
A random vector $\state \in \R^n$ is \textbf{Gaussian}\index{Gaussian!multivariate} with mean $\vect{\mu}$ and covariance $\mat{\Sigma} \succ \mat{0}$, written $\state \sim \Normal(\vect{\mu}, \mat{\Sigma})$, if its density is
\begin{equation}
  \Normal(\state; \vect{\mu}, \mat{\Sigma}) = \frac{1}{(2\pi)^{n/2}\sqrt{\det\mat{\Sigma}}}\exp\left(-\tfrac12\,d^2(\state)\right),\qquad
  d^2(\state) = (\state - \vect{\mu})\T\mat{\Sigma}\inv(\state - \vect{\mu}),
  \label{eq:appB-gaussian}
\end{equation}
where $d(\state)$ is the \textbf{Mahalanobis distance}\index{Mahalanobis distance} of $\state$ from $\vect{\mu}$. Then $\E[\state] = \vect{\mu}$ and $\Cov(\state) = \mat{\Sigma}$.
\end{definition}

The density is constant on the ellipsoids $d^2(\state) = k^2$. By
\cref{thm:appB-spectral} applied to $\mat{\Sigma} =
\mat{Q}\mat{\Lambda}\mat{Q}\T$, the ellipsoid $\mathcal{E}_k$ of \cref{ch:ch02}
is centred at $\vect{\mu}$, its axes are the eigenvectors $\vect{q}_i$ and its
semi-axes are $k\sqrt{\lambda_i}$: a large eigenvalue is a direction of large
uncertainty. In two dimensions the probability mass inside $\mathcal{E}_k$ is
$1 - e^{-k^2/2}$, that is $39.3\,\%$ for $k = 1$ and $86.5\,\%$ for $k = 2$, a
consequence of \cref{thm:appB-chi2} below. A singular $\mat{\Sigma} \succeq
\mat{0}$ has no density, but the two theorems that follow still hold for it.

\begin{theorem}[Affine transformations of Gaussians]
\label{thm:appB-affine}
If $\state \sim \Normal(\vect{\mu}, \mat{\Sigma})$ in $\R^n$ and $\vect{y} = \mat{A}\state + \vect{b}$ with $\mat{A} \in \R^{m \times n}$, then $\vect{y} \sim \Normal(\mat{A}\vect{\mu} + \vect{b},\ \mat{A}\mat{\Sigma}\mat{A}\T)$.
\end{theorem}
\begin{proof}
The mean and the covariance follow from \cref{thm:appB-cov-rules}. For the shape, take first $m = n$ and $\mat{A}$ invertible: substituting $\state = \mat{A}\inv(\vect{y} - \vect{b})$ into \cref{eq:appB-gaussian} and multiplying by the change-of-variables factor $\abs{\det\mat{A}\inv}$ gives exactly the density $\Normal(\vect{y}; \mat{A}\vect{\mu} + \vect{b}, \mat{A}\mat{\Sigma}\mat{A}\T)$, because $(\mat{A}\mat{\Sigma}\mat{A}\T)\inv = \mat{A}^{-\mathsf{T}}\mat{\Sigma}\inv\mat{A}\inv$ and $\det(\mat{A}\mat{\Sigma}\mat{A}\T) = (\det\mat{A})^2\det\mat{\Sigma}$. A non-square $\mat{A}$ of full row rank is extended to an invertible matrix, and the extra coordinates are integrated out with the marginalisation half of \cref{thm:appB-conditioning}, whose proof does not use this theorem; the remaining cases follow with characteristic functions \cite[Chapter~4]{sarkka2013bayesian}. \qedhere
\end{proof}

In particular the sum of independent Gaussians is Gaussian with the sums of
the means and of the covariances (apply $\mat{A} = [\mat{I}\ \ \mat{I}]$ to
the joint vector, whose covariance is block-diagonal), which is how process
noise enters the prediction step of \cref{ch:ch18}.

\begin{theorem}[Marginal and conditional of a joint Gaussian]
\label{thm:appB-conditioning}
Let $\state = (\state_a, \state_b)$ be jointly Gaussian,
\begin{equation}
  \begin{bmatrix}\state_a\\ \state_b\end{bmatrix} \sim
  \Normal\left(\begin{bmatrix}\vect{\mu}_a\\ \vect{\mu}_b\end{bmatrix},
  \begin{bmatrix}\mat{\Sigma}_{aa} & \mat{\Sigma}_{ab}\\ \mat{\Sigma}_{ba} & \mat{\Sigma}_{bb}\end{bmatrix}\right),
  \qquad \mat{\Sigma}_{ba} = \mat{\Sigma}_{ab}\T .
  \label{eq:appB-joint}
\end{equation}
Then the \textbf{marginal}\index{Gaussian!marginal} of $\state_b$ is $\Normal(\vect{\mu}_b, \mat{\Sigma}_{bb})$ (and likewise for $\state_a$), and the \textbf{conditional}\index{Gaussian!conditional} density of $\state_a$ given $\state_b$ is Gaussian,
\begin{equation}
  \state_a \mid \state_b \sim \Normal\left(\vect{\mu}_a + \mat{\Sigma}_{ab}\mat{\Sigma}_{bb}\inv(\state_b - \vect{\mu}_b),\;
  \mat{\Sigma}_{aa} - \mat{\Sigma}_{ab}\mat{\Sigma}_{bb}\inv\mat{\Sigma}_{ba}\right).
  \label{eq:appB-conditional}
\end{equation}
The conditional mean is affine in the observed value; the conditional covariance is the Schur complement of $\mat{\Sigma}_{bb}$, does not depend on the observed value, and is never larger than $\mat{\Sigma}_{aa}$.
\end{theorem}
\begin{proof}
Write $\vect{d}_a = \state_a - \vect{\mu}_a$, $\vect{d}_b = \state_b - \vect{\mu}_b$, $\mat{S} = \mat{\Sigma}_{aa} - \mat{\Sigma}_{ab}\mat{\Sigma}_{bb}\inv\mat{\Sigma}_{ba}$ and $\mat{\Lambda} = \mat{\Sigma}\inv$ with blocks $\mat{\Lambda}_{aa}, \mat{\Lambda}_{ab}, \mat{\Lambda}_{ba}, \mat{\Lambda}_{bb}$. By \cref{thm:appB-block} with $\mat{D} = \mat{\Sigma}_{bb}$,
\begin{equation*}
  \mat{\Lambda}_{aa} = \mat{S}\inv,\qquad
  \mat{\Lambda}_{ab} = -\mat{S}\inv\mat{\Sigma}_{ab}\mat{\Sigma}_{bb}\inv,\qquad
  \det\mat{\Sigma} = \det\mat{S}\,\det\mat{\Sigma}_{bb} .
\end{equation*}
The exponent of the joint density is $-\tfrac12(\vect{d}_a\T\mat{\Lambda}_{aa}\vect{d}_a + 2\vect{d}_a\T\mat{\Lambda}_{ab}\vect{d}_b + \vect{d}_b\T\mat{\Lambda}_{bb}\vect{d}_b)$. Completing the square in $\vect{d}_a$ with $\vect{m} = -\mat{\Lambda}_{aa}\inv\mat{\Lambda}_{ab}\vect{d}_b = \mat{\Sigma}_{ab}\mat{\Sigma}_{bb}\inv\vect{d}_b$ turns it into
\begin{equation*}
  -\tfrac12(\vect{d}_a - \vect{m})\T\mat{S}\inv(\vect{d}_a - \vect{m})
  -\tfrac12\vect{d}_b\T\left(\mat{\Lambda}_{bb} - \mat{\Lambda}_{ba}\mat{\Lambda}_{aa}\inv\mat{\Lambda}_{ab}\right)\vect{d}_b ,
\end{equation*}
and the matrix in the second term is the Schur complement of $\mat{\Lambda}_{aa}$ in $\mat{\Lambda}$, which by \cref{thm:appB-block} applied to $\mat{\Lambda}$ (whose inverse is $\mat{\Sigma}$) equals $\mat{\Sigma}_{bb}\inv$. Splitting the normalising constant with $\det\mat{\Sigma} = \det\mat{S}\det\mat{\Sigma}_{bb}$ and $(2\pi)^{n} = (2\pi)^{n_a}(2\pi)^{n_b}$ gives the exact factorisation
\begin{equation}
  p(\state_a, \state_b) = \Normal\left(\state_a;\ \vect{\mu}_a + \vect{m},\ \mat{S}\right)\cdot \Normal\left(\state_b;\ \vect{\mu}_b,\ \mat{\Sigma}_{bb}\right).
  \label{eq:appB-factorisation}
\end{equation}
Integrating over $\state_a$ removes the first factor, which is a density in $\state_a$, and leaves the marginal of $\state_b$. Dividing by that marginal gives $p(\state_a \mid \state_b) = p(\state_a, \state_b)/p(\state_b)$, the first factor, which is \cref{eq:appB-conditional}. Finally $\mat{\Sigma}_{aa} - \mat{S} = \mat{\Sigma}_{ab}\mat{\Sigma}_{bb}\inv\mat{\Sigma}_{ba} \succeq \mat{0}$ by \cref{thm:appB-psd-tests}(iii). \qedhere
\end{proof}

\begin{corollary}[Linear-Gaussian measurement update]
\label{thm:appB-linear-update}
Let $\state \sim \Normal(\vect{\mu}, \mat{P})$ and $\meas = \mat{H}\state + \vect{v}$ with $\vect{v} \sim \Normal(\vect{0}, \mat{R})$ independent of $\state$. Then $(\state, \meas)$ is jointly Gaussian with $\Cov(\meas) = \mat{H}\mat{P}\mat{H}\T + \mat{R} =: \mat{S}$ and $\Cov(\state, \meas) = \mat{P}\mat{H}\T$, and
\begin{equation}
  \state \mid \meas \sim \Normal\left(\vect{\mu} + \Kgain(\meas - \mat{H}\vect{\mu}),\ \mat{P} - \Kgain\mat{H}\mat{P}\right),\qquad \Kgain = \mat{P}\mat{H}\T\mat{S}\inv .
  \label{eq:appB-kalman-update}
\end{equation}
\end{corollary}
\begin{proof}
$(\state, \meas)$ is the image of the Gaussian vector $(\state, \vect{v})$, whose covariance is block-diagonal, under the linear map $\begin{bsmallmatrix}\mat{I} & \mat{0}\\ \mat{H} & \mat{I}\end{bsmallmatrix}$, hence Gaussian by \cref{thm:appB-affine}; its blocks follow from \cref{thm:appB-cov-rules}. Insert $\mat{\Sigma}_{ab} = \mat{P}\mat{H}\T$ and $\mat{\Sigma}_{bb} = \mat{S}$ into \cref{eq:appB-conditional}. \qedhere
\end{proof}

\Cref{eq:appB-kalman-update} is the update step of the Kalman filter with
gain $\Kgain$: \cref{ch:ch18} derives it from this corollary, and
\cref{thm:appB-woodbury} turns $\mat{P} - \Kgain\mat{H}\mat{P}$ into the
information form $(\mat{P}\inv + \mat{H}\T\mat{R}\inv\mat{H})\inv$, in which
the precisions of prior and measurement simply add.

\begin{example}[Conditioning the Gaussian of \cref{ex:appB-eigen}]
\label{ex:appB-conditioning}
Let $(x, y) \sim \Normal(\vect{\mu}, \mat{\Sigma})$ with $\vect{\mu} = (2, 1)\T$ and $\mat{\Sigma}$ from \cref{eq:appB-sigma}, and suppose $y = 2$ is observed. \Cref{eq:appB-conditional} with $a = x$ and $b = y$ gives
\begin{equation*}
  \E[x \mid y = 2] = 2 + \frac{1.2}{1.0}(2 - 1) = 3.2,\qquad
  \Var(x \mid y = 2) = 2.0 - \frac{1.2^2}{1.0} = 0.56 .
\end{equation*}
Before the observation $x \sim \Normal(2, 2.0)$, with standard deviation $1.414$; afterwards $x \mid y \sim \Normal(3.2, 0.56)$, with standard deviation $0.748$. The positive correlation $\rho = 1.2/\sqrt{2.0 \cdot 1.0} = 0.85$ means that a $y$ above its mean shifts the belief about $x$ upwards, and the variance shrinks by the factor $1 - \rho^2 = 0.28$. \Cref{fig:appB-conditioning} shows the ellipses, the slice $y = 2$ and both densities.
\end{example}

\begin{figure}[tb]
  \centering
  \inputfigure{appB/gaussian-conditioning}
  \caption{Conditioning the Gaussian of \cref{ex:appB-conditioning}. The $1\sigma$ and $2\sigma$ ellipses of the joint density are tilted along the eigenvector $\vect{q}_1$ of $\mat{\Sigma}$. Observing $y = 2$ restricts the density to the dashed slice; along it, the conditional density of $x$ (orange) is centred at $3.2$, to the right of the prior mean $2$, and is narrower than the marginal density of $x$ (blue, drawn along the bottom) because knowing $y$ removes the part of the uncertainty of $x$ that is correlated with $y$.}
  \label{fig:appB-conditioning}
\end{figure}

\begin{table}[tb]
  \centering\small
  \caption{Gaussian identities. $\state \sim \Normal(\vect{\mu}, \mat{\Sigma})$, blocks as in \cref{eq:appB-joint}; ``$\propto$'' means equality up to a factor that does not depend on $\state$.}
  \label{tab:appB-gaussian}
  \begin{tabular}{@{}L{9.0cm}L{5.8cm}@{}}
  \toprule
  Identity & Remark\\
  \midrule
  $\mat{A}\state + \vect{b} \sim \Normal(\mat{A}\vect{\mu} + \vect{b},\ \mat{A}\mat{\Sigma}\mat{A}\T)$ & \cref{thm:appB-affine}; prediction step\\
  $\state_a \sim \Normal(\vect{\mu}_a, \mat{\Sigma}_{aa})$ & marginal: drop the other blocks\\
  $\state_a \mid \state_b \sim \Normal(\vect{\mu}_{a|b}, \mat{\Sigma}_{a|b})$, \ $\vect{\mu}_{a|b} = \vect{\mu}_a + \mat{\Sigma}_{ab}\mat{\Sigma}_{bb}\inv(\state_b - \vect{\mu}_b)$, \ $\mat{\Sigma}_{a|b} = \mat{\Sigma}_{aa} - \mat{\Sigma}_{ab}\mat{\Sigma}_{bb}\inv\mat{\Sigma}_{ba}$ & \cref{thm:appB-conditioning}; update step\\
  $\state + \vect{y} \sim \Normal(\vect{\mu}_x + \vect{\mu}_y,\ \mat{\Sigma}_x + \mat{\Sigma}_y)$ & $\state$, $\vect{y}$ independent\\
  $\Normal(\state; \vect{a}, \mat{A})\,\Normal(\state; \vect{b}, \mat{B}) \propto \Normal(\state; \vect{c}, \mat{C})$, \ $\mat{C} = (\mat{A}\inv + \mat{B}\inv)\inv$, \ $\vect{c} = \mat{C}(\mat{A}\inv\vect{a} + \mat{B}\inv\vect{b})$ & precisions add; information form of the update\\
  $d^2(\state) = (\state - \vect{\mu})\T\mat{\Sigma}\inv(\state - \vect{\mu}) \sim \chi^2_n$ & \cref{thm:appB-chi2}; gating\\
  $\state = \vect{\mu} + \mat{L}\vect{\xi}$, \ $\mat{L}\mat{L}\T = \mat{\Sigma}$, \ $\vect{\xi} \sim \Normal(\vect{0}, \mat{I})$ & how to draw a sample\\
  \bottomrule
  \end{tabular}
\end{table}

\subsection{Bayes' rule, total probability and the recursive filter}

For two random vectors the \textbf{chain rule} $p(\state, \meas) = p(\meas
\mid \state)\,p(\state)$ defines the conditional density, and integrating it
over $\state$ gives the law of \textbf{total probability}\index{total
probability}, $p(\meas) = \int p(\meas \mid \state)p(\state)\,\mathrm{d}\state$.
Reading the chain rule in both directions gives \textbf{Bayes'
rule}\index{Bayes' rule}
\begin{equation}
  p(\state \mid \meas) = \frac{p(\meas \mid \state)\,p(\state)}{p(\meas)}
  = \frac{p(\meas \mid \state)\,p(\state)}{\int p(\meas \mid \state')\,p(\state')\,\mathrm{d}\state'} ,
  \label{eq:appB-bayes}
\end{equation}
which turns a \textbf{prior} $p(\state)$ and a \textbf{likelihood} $p(\meas
\mid \state)$ into the \textbf{posterior} $p(\state \mid \meas)$; the
denominator is a constant that makes the posterior integrate to one. For
discrete events the integrals are sums. A small example: a detector flags an
intruder with probability $0.9$ when one is present and with probability
$0.2$ when none is (a false alarm); if intruders are present $10\,\%$ of the
time, a flag raises the probability of an intruder only to $0.09/(0.09 +
0.18) = 1/3$, because false alarms on the common ``no intruder'' case
outnumber true detections on the rare one.

Every filter of \cref{ch:ch18,ch:ch19} is these two rules applied over and
over: with the Markov assumption $p(\state_k \mid \state_{k-1},
\meas_{1:k-1}) = p(\state_k \mid \state_{k-1})$ and measurements that depend
only on the current state, the belief about $\state_k$ given $\meas_{1:k}$ is
obtained in two steps,
\begin{align}
  p(\state_k \mid \meas_{1:k-1}) &= \int p(\state_k \mid \state_{k-1})\,p(\state_{k-1} \mid \meas_{1:k-1})\,\mathrm{d}\state_{k-1} && \text{(predict: total probability)},
  \label{eq:appB-predict}\\
  p(\state_k \mid \meas_{1:k}) &\propto p(\meas_k \mid \state_k)\,p(\state_k \mid \meas_{1:k-1}) && \text{(update: Bayes' rule)}.
  \label{eq:appB-update}
\end{align}
When the motion and measurement models are linear with Gaussian noise, both
steps map Gaussians to Gaussians (\cref{thm:appB-affine,thm:appB-linear-update})
and the recursion is the Kalman filter; when they are not, \cref{ch:ch19}
approximates the two integrals by linearisation, by sigma points or by
samples \cite{thrun2005probabilistic,sarkka2013bayesian}.

\subsection{Monte Carlo estimation}
\label{sec:appB-montecarlo}

Many quantities in this book are expectations without a closed form, such
as the probability that a predicted trajectory hits an obstacle or the
collision rate of a planner over random scenarios. The \textbf{Monte
Carlo}\index{Monte Carlo estimation} estimate of $\E[g(\state)]$ from $N$
independent samples $\state^{(1)}, \dots, \state^{(N)}$ drawn from $p$ is
\begin{equation}
  \hat{g}_N = \frac1N\sum_{i=1}^{N} g(\state^{(i)}),\qquad
  \E[\hat{g}_N] = \E[g(\state)],\qquad
  \Var(\hat{g}_N) = \frac{\Var(g(\state))}{N} .
  \label{eq:appB-montecarlo}
\end{equation}
The estimate is unbiased and its standard error shrinks like $1/\sqrt{N}$
whatever the dimension of $\state$. For a probability $\Prob(\state \in
\mathcal{R})$ take $g$ the indicator of $\mathcal{R}$; then the standard error
is $\sqrt{\hat{p}(1-\hat{p})/N}$, the error bar to attach to a collision rate
estimated from $N$ random scenarios (\cref{ch:ch25}). To draw from
$\Normal(\vect{\mu}, \mat{\Sigma})$, draw $\vect{\xi} \sim \Normal(\vect{0},
\mat{I})$ and set $\state = \vect{\mu} + \mat{L}\vect{\xi}$ with the Cholesky
factor $\mat{L}\mat{L}\T = \mat{\Sigma}$ (\cref{thm:appB-affine} in reverse). Drawing $N = 10\,000$ samples this way from the Gaussian of \cref{ex:appB-conditioning} and counting those with $d^2 \le 4$ gives $\hat{p} = 0.8695$ with standard error $0.0034$, within $1.5$ standard errors of the exact value $1 - e^{-2} = 0.8647$ of \cref{thm:appB-chi2} below.
When samples can only be drawn from another density $q$, \textbf{importance
sampling}\index{importance sampling} weights them by $w^{(i)} \propto
p(\state^{(i)})/q(\state^{(i)})$, normalised to sum to one; the particle
filter of \cref{ch:ch19} is a sequential version of this idea
\cite{doucet2001sequential}.

\subsection{Mahalanobis distance and chi-square gating}
\label{sec:appB-gating}

\begin{theorem}[Mahalanobis distance of a Gaussian vector]
\label{thm:appB-chi2}
If $\state \sim \Normal(\vect{\mu}, \mat{\Sigma})$ in $\R^n$ with $\mat{\Sigma} \succ \mat{0}$, then $d^2(\state) = (\state - \vect{\mu})\T\mat{\Sigma}\inv(\state - \vect{\mu})$ has the \textbf{chi-square distribution}\index{chi-square distribution} with $n$ degrees of freedom, $\chi^2_n$, whose mean is $n$. For $n = 2$, $\Prob(d^2 \le \gamma) = 1 - e^{-\gamma/2}$.
\end{theorem}
\begin{proof}
Let $\mat{L}\mat{L}\T = \mat{\Sigma}$ and $\vect{\xi} = \mat{L}\inv(\state - \vect{\mu})$. By \cref{thm:appB-affine}, $\vect{\xi} \sim \Normal(\vect{0}, \mat{L}\inv\mat{\Sigma}\mat{L}^{-\mathsf{T}}) = \Normal(\vect{0}, \mat{I})$, so its entries are $n$ independent standard normal variables, and $d^2 = \vect{\xi}\T\vect{\xi} = \sum_i \xi_i^2$ is by definition $\chi^2_n$ distributed; $\E[\xi_i^2] = 1$ gives the mean. For $n = 2$ the density of $(\xi_1, \xi_2)$ is $\frac{1}{2\pi}e^{-\varrho^2/2}$ in polar coordinates $(\varrho, \phi)$, and $\Prob(d^2 \le \gamma) = \int_0^{\sqrt{\gamma}} e^{-\varrho^2/2}\varrho\,\mathrm{d}\varrho = 1 - e^{-\gamma/2}$. \qedhere
\end{proof}

The theorem is the basis of \textbf{gating}\index{gating}: a filter that
expects a measurement with innovation $\vect{y} = \meas - \hat{\meas}$ of
covariance $\mat{S}$ accepts it only if $\vect{y}\T\mat{S}\inv\vect{y} \le
\gamma$, with $\gamma$ the quantile of $\chi^2_n$ at the chosen confidence $p$:
for $n = 2$ the $95\,\%$ and $99\,\%$ gates are $\gamma = -2\ln(1-p) = 5.991$
and $9.210$, for $n = 1$ they are $3.841$ and $6.635$, and for $n = 3$ they
are $7.815$ and $11.345$; a true measurement is then wrongly rejected with
probability $1 - p$, and a detection from another target or from clutter is
rejected whenever it lands outside the ellipsoid
\cite{barshalom2001estimation}. This is the gate of \cref{ch:ch18,ch:ch19},
and the same quantiles size the uncertainty-inflated constraints of
\cref{ch:ch24}, where an obstacle known only as a Gaussian is enlarged to the
ellipse that holds a required probability mass.

% ---------------------------------------------------------------------
\section{Optimization}
\label{sec:appB-optimization}

\subsection{Convex sets and convex functions}

\begin{definition}[Convex set, convex function]
\label{def:appB-convex}
A set $\mathcal{C} \subseteq \R^n$ is \textbf{convex}\index{convexity!set} if the segment between any two of its points stays inside it: $\vect{x}, \vect{y} \in \mathcal{C}$ and $\theta \in [0, 1]$ imply $\theta\vect{x} + (1 - \theta)\vect{y} \in \mathcal{C}$. A function $f\colon \mathcal{C} \to \R$ on a convex set is \textbf{convex}\index{convexity!function} if its graph lies below every chord,
\begin{equation}
  f(\theta\vect{x} + (1 - \theta)\vect{y}) \le \theta f(\vect{x}) + (1 - \theta) f(\vect{y})\qquad \text{for all } \vect{x}, \vect{y} \in \mathcal{C},\ \theta \in [0, 1],
  \label{eq:appB-convex}
\end{equation}
strictly convex if the inequality is strict for $\vect{x} \ne \vect{y}$ and $\theta \in (0, 1)$, and concave if $-f$ is convex.
\end{definition}

Half-spaces, balls, ellipsoids and affine subspaces are convex, and any
intersection of convex sets is convex: the velocities allowed by all \orca
half-planes in \cref{ch:ch13} form a convex polygon, whereas the
collision-free velocities outside a velocity obstacle of \cref{ch:ch12} do not
form a convex set (\cref{fig:appB-convex}). Linear functions, norms,
$\norm{\vect{x} - \vect{c}}^2$ and $\vect{x}\T\mat{A}\vect{x}$ with $\mat{A}
\succeq \mat{0}$ are convex, as are non-negative sums and pointwise maxima of
convex functions, and the sublevel sets $\set{\vect{x} : f(\vect{x}) \le
\alpha}$ of a convex function are convex. What makes convexity the dividing
line of optimization is the following fact.

\begin{proposition}[Local minima of convex functions are global]
\label{thm:appB-local-global}
Let $f$ be convex on a convex set $\mathcal{C}$ and let $\vect{x}^*$ be a local minimiser: $f(\vect{x}^*) \le f(\vect{x})$ for all $\vect{x} \in \mathcal{C}$ within some distance $\eps > 0$ of $\vect{x}^*$. Then $f(\vect{x}^*) \le f(\vect{x})$ for all $\vect{x} \in \mathcal{C}$. If $f$ is differentiable on $\R^n$, $\vect{x}^*$ is a global minimiser iff $\nabla f(\vect{x}^*) = \vect{0}$.
\end{proposition}
\begin{proof}
Suppose $f(\vect{y}) < f(\vect{x}^*)$ for some $\vect{y} \in \mathcal{C}$. The points $\vect{x}_\theta = \theta\vect{y} + (1-\theta)\vect{x}^*$ lie in $\mathcal{C}$ and satisfy $f(\vect{x}_\theta) \le \theta f(\vect{y}) + (1-\theta)f(\vect{x}^*) < f(\vect{x}^*)$ for every $\theta \in (0, 1]$; for $\theta$ small enough $\vect{x}_\theta$ is within $\eps$ of $\vect{x}^*$, contradicting local minimality. For the second claim, a differentiable convex function satisfies $f(\vect{y}) \ge f(\vect{x}) + \nabla f(\vect{x})\T(\vect{y} - \vect{x})$ for all $\vect{x}, \vect{y}$, that is, its graph lies above its tangent planes \cite[Section~3.1]{boyd2004convex}, so a zero gradient makes $\vect{x}$ a global minimiser; conversely, a minimiser of a differentiable function has zero gradient. \qedhere
\end{proof}

For twice differentiable $f$ the test is the Hessian: $f$ is convex on $\R^n$
iff $\nabla^2 f(\vect{x}) \succeq \mat{0}$ everywhere
\cite[Section~3.1]{boyd2004convex}. This is why the quadratic costs of
\cref{ch:ch21} are built from PSD weights, and why the potential fields of
\cref{ch:ch15}, whose repulsive terms are not convex, can trap a drone in a
local minimum.

\begin{figure}[tb]
  \centering
  \inputfigure{appB/convexity}
  \caption{Convexity. Left: in a convex set every chord stays inside; in the non-convex set the dashed chord leaves it. Middle: an intersection of half-planes, the situation of \orca in \cref{ch:ch13}, is convex, while the set of velocities outside a cone, the situation of a velocity obstacle in \cref{ch:ch12}, is not. Right: a convex function lies below its chords and above its tangents, so its only stationary point is the global minimum; the non-convex function has a local minimum in which gradient descent can stop.}
  \label{fig:appB-convex}
\end{figure}

\subsection{Optimality conditions: Lagrange multipliers and KKT}
\label{sec:appB-kkt}

The general constrained problem is
\begin{equation}
  \min_{\vect{x} \in \R^n}\ f(\vect{x})\quad \st g_i(\vect{x}) \le 0\ \ (i = 1, \dots, m),\qquad h_j(\vect{x}) = 0\ \ (j = 1, \dots, p),
  \label{eq:appB-nlp}
\end{equation}
with differentiable $f$, $g_i$ and $h_j$. Its
\textbf{Lagrangian}\index{Lagrangian} attaches a multiplier to every
constraint, $\mathcal{L}(\vect{x}, \vect{\lambda}, \vect{\nu}) = f(\vect{x}) +
\sum_i \lambda_i g_i(\vect{x}) + \sum_j \nu_j h_j(\vect{x})$.

\begin{theorem}[Karush--Kuhn--Tucker conditions]
\label{thm:appB-kkt}
If $\vect{x}^*$ is a local minimiser of \cref{eq:appB-nlp} and the constraints satisfy a constraint qualification at $\vect{x}^*$ (for instance, the gradients of the active constraints are linearly independent, or all constraints are affine), then there exist multipliers $\vect{\lambda}^* \in \R^m$ and $\vect{\nu}^* \in \R^p$ with\index{KKT conditions}
\begin{align}
  \nabla f(\vect{x}^*) + \sum_i \lambda_i^*\nabla g_i(\vect{x}^*) + \sum_j \nu_j^*\nabla h_j(\vect{x}^*) &= \vect{0} && \text{(stationarity)},\label{eq:appB-kkt-stat}\\
  g_i(\vect{x}^*) \le 0,\quad h_j(\vect{x}^*) &= 0 && \text{(primal feasibility)},\\
  \lambda_i^* &\ge 0 && \text{(dual feasibility)},\\
  \lambda_i^*\,g_i(\vect{x}^*) &= 0 && \text{(complementary slackness)}.\label{eq:appB-kkt-cs}
\end{align}
If $f$ and the $g_i$ are convex and the $h_j$ affine, these conditions are also sufficient: every point that satisfies them is a global minimiser.
\end{theorem}

The proof is in \textcite[Chapter~12]{nocedal2006numerical} and, for the
convex case, in \textcite[Chapter~5]{boyd2004convex}. The conditions say that
at the optimum the negative gradient of the cost is balanced by the normals of
the \emph{active} constraints, those with $g_i(\vect{x}^*) = 0$, which by
\cref{eq:appB-kkt-cs} are the only ones that can carry a positive multiplier;
$\lambda_i^*$ measures how much the optimal cost would drop if constraint $i$
were relaxed by one unit. This is the theory behind the soft constraints of
\cref{ch:ch21}: a slack variable with a linear penalty reproduces the hard
constraint exactly, whenever it is feasible, as soon as the penalty weight
exceeds the multiplier \cite[Chapter~17]{nocedal2006numerical}. With equality
constraints only, the conditions reduce to the classical Lagrange multipliers.

\subsection{Linear programs}
\label{sec:appB-lp}

\begin{definition}[Linear program]
\label{def:appB-lp}
A \textbf{linear program}\index{linear program} (LP) minimises a linear function over a polyhedron:
\begin{equation}
  \min_{\vect{z} \in \R^n}\ \vect{c}\T\vect{z}\quad \st \mat{A}\vect{z} \le \vect{b},\qquad P = \set{\vect{z} : \mat{A}\vect{z} \le \vect{b}} \text{ the feasible polyhedron}.
  \label{eq:appB-lp}
\end{equation}
Equality constraints, sign constraints $\vect{z} \ge \vect{0}$ and maximisation are all reduced to this form by writing $\vect{a}\T\vect{z} = b$ as two inequalities, $z_j \ge 0$ as $-z_j \le 0$ and $\max \vect{c}\T\vect{z}$ as $\min (-\vect{c})\T\vect{z}$. Simplex solvers and most textbooks use instead the \textbf{standard form}\index{linear program!standard form} $\min \vect{c}\T\vect{z}$ subject to $\mat{A}\vect{z} = \vect{b}$, $\vect{z} \ge \vect{0}$, reached from \cref{eq:appB-lp} by adding a \emph{slack} variable $s_i \ge 0$ to every inequality, $\vect{a}_i\T\vect{z} + s_i = b_i$, and by splitting a free variable into $z_j = z_j^+ - z_j^-$ with $z_j^\pm \ge 0$; \code{scipy.optimize.linprog} accepts inequalities, equalities and bounds directly.
\end{definition}

An LP is convex, so \cref{thm:appB-kkt} applies with $\nabla g_i =
\vect{a}_i$: at an optimum, $-\vect{c}$ is a non-negative combination of the
normals of the active constraints. In two dimensions this is a picture
(\cref{fig:appB-lp}): $P$ is a polygon, the level lines $\vect{c}\T\vect{z} =
\text{const}$ are parallel, and sliding a level line in the direction
$-\vect{c}$ as far as $P$ allows ends at a vertex, or along a whole edge when
the level lines are parallel to it.

\begin{theorem}[Vertex optimality]
\label{thm:appB-vertex}
If the LP \cref{eq:appB-lp} has an optimal solution and $P$ has at least one vertex, then some vertex of $P$ is optimal. Otherwise the LP is infeasible ($P = \emptyset$) or unbounded ($\vect{c}\T\vect{z} \to -\infty$ on $P$).
\end{theorem}

A proof is in \textcite[Chapter~13]{nocedal2006numerical}. The simplex
method walks from vertex to adjacent vertex, improving the objective at every
step until no neighbour is better, which by convexity is a global optimum;
interior-point methods approach the optimum through the inside of $P$ and are
polynomial in the worst case. \Cref{ch:ch22} calls such solvers through
\code{scipy.optimize}, and branch and bound (\cref{sec:appB-milp}) solves a
sequence of LPs. The velocity choice of \orca in \cref{ch:ch13} is the same
geometry: the velocity closest to the preferred one inside an intersection of
half-planes, found by a small linear-programming routine that adds one
half-plane at a time.

\begin{example}[A two-variable LP]
\label{ex:appB-lp}
Maximise $2z_1 + 3z_2$ subject to $z_1 + z_2 \le 4$, $z_1 + 3z_2 \le 6$ and $z_1, z_2 \ge 0$. The polygon $P$ has the vertices $(0,0)$, $(4,0)$, $(3,1)$ and $(0,2)$, where the objective takes the values $0$, $8$, $9$ and $6$; the optimum is $\vect{z}^* = (3, 1)$ with value $9$ (\cref{fig:appB-lp}). At $\vect{z}^*$ both structural constraints are active, and the stationarity condition for the maximisation reads $\vect{c} = \lambda_1(1, 1)\T + \lambda_2(1, 3)\T$, solved by $\lambda_1 = 1.5$ and $\lambda_2 = 0.5$: both non-negative, as required, so $\vect{c}$ lies in the cone spanned by the two active normals and no feasible direction improves the objective. Relaxing the first constraint to $z_1 + z_2 \le 5$ raises the optimum by exactly $\lambda_1 = 1.5$, to $10.5$ at $(4.5, 0.5)$.
\end{example}

\begin{figure}[tb]
  \centering
  \inputfigure{appB/lp-geometry}
  \caption{The LP of \cref{ex:appB-lp}. The shaded polygon is the feasible region $P$, cut out by four half-planes; the dotted lines are level lines of the objective $2z_1 + 3z_2$. Sliding a level line in the direction of $\vect{c}$ until it is about to leave $P$ stops at the vertex $(3, 1)$, where $\vect{c}$ lies between the normals of the two active constraints.}
  \label{fig:appB-lp}
\end{figure}

\subsection{Quadratic programs}
\label{sec:appB-qp}

A \textbf{quadratic program}\index{quadratic program} (QP) has a quadratic
objective and linear constraints,
\begin{equation}
  \min_{\vect{z}}\ \tfrac12\vect{z}\T\mat{H}\vect{z} + \vect{f}\T\vect{z}\quad \st \vect{l} \le \mat{G}\vect{z} \le \vect{u},
  \label{eq:appB-qp}
\end{equation}
which is the form of the MPC problem of \cref{ch:ch21}. It is convex iff
$\mat{H} \succeq \mat{0}$, and then \cref{thm:appB-kkt} is necessary and
sufficient. Without constraints the solution is the linear system
$\mat{H}\vect{z} = -\vect{f}$, unique when $\mat{H} \succ \mat{0}$; the
least-squares problem of \cref{thm:appB-normal} is the case $\mat{H} =
2\mat{A}\T\mat{A}$, $\vect{f} = -2\mat{A}\T\vect{b}$. With equality
constraints $\mat{G}\vect{z} = \vect{g}$ only, the KKT conditions are the
single linear system
\begin{equation}
  \begin{bmatrix}\mat{H} & \mat{G}\T\\ \mat{G} & \mat{0}\end{bmatrix}
  \begin{bmatrix}\vect{z}\\ \vect{\nu}\end{bmatrix} =
  \begin{bmatrix}-\vect{f}\\ \vect{g}\end{bmatrix},
  \label{eq:appB-kkt-system}
\end{equation}
a block matrix whose Schur complement $-\mat{G}\mat{H}\inv\mat{G}\T$
(\cref{thm:appB-block}) is what a solver eliminates first. Inequality
constraints are handled by active-set methods, which guess the active set and
solve \cref{eq:appB-kkt-system} for it, by interior-point methods
\cite{nocedal2006numerical}, or by operator-splitting methods such as ADMM,
which alternate projections onto the constraints with linear solves.

\subsection{Mixed-integer programs, branch and bound, big-M}
\label{sec:appB-milp}

A \textbf{mixed-integer linear program}\index{mixed-integer linear program}
(MILP) is an LP in which some variables are restricted to integers, most
often to the binary values $b \in \set{0, 1}$. The integrality constraint
destroys convexity: the feasible set is a scatter of points, not a
polyhedron, and the problem is NP-hard in general. Dropping the integrality
gives the \textbf{LP relaxation}\index{LP relaxation}, whose optimal value is
a bound on the MILP optimum (a lower bound for a minimisation, an upper
bound for a maximisation), because the relaxation optimises over a larger
set. \textbf{Branch and bound}\index{branch and bound} turns this bound into
an exact method \cite{wolsey1998integer}:
\begin{enumerate}
  \item Solve the LP relaxation of the current subproblem. If it is infeasible, or if its bound cannot beat the best integer solution found so far (the \emph{incumbent}), discard the subproblem (\emph{prune}).
  \item If the relaxed solution is integral, it is the best solution of this subproblem; make it the incumbent if it is better than the current one.
  \item Otherwise pick a variable $z_j$ with fractional value $v$ and \emph{branch}: create two subproblems with the extra constraint $z_j \le \lfloor v\rfloor$ and $z_j \ge \lceil v\rceil$ respectively, and repeat.
\end{enumerate}
The method terminates because every branch cuts away a slice of the relaxed
region that contains no integer point, and it is exact because a subproblem is
pruned only when it provably contains nothing better than the incumbent; its
running time is exponential in the worst case but often small in practice,
the trade-off that \cref{ch:ch22} explores.

\begin{example}[Branch and bound by hand]
\label{ex:appB-bnb}
Maximise $5z_1 + 4z_2$ subject to $6z_1 + 4z_2 \le 24$, $z_1 + 2z_2 \le 6$, $z_1, z_2 \ge 0$ and integer. The relaxation (node~1) has the optimum $(3, 1.5)$ with value $21$, so no integer solution can exceed $21$. Branching on $z_2$ gives node~2 ($z_2 \le 1$) with relaxed optimum $(3.33, 1)$ and bound $20.67$, and node~3 ($z_2 \ge 2$) with $(2, 2)$ and bound $18$. Branching node~2 on $z_1$ gives node~4 ($z_1 \le 3$), whose relaxation $(3, 1)$ is integral with value $19$ and becomes the incumbent, and node~5 ($z_1 \ge 4$), whose relaxation $(4, 0)$ is integral with value $20$ and replaces it. Node~3 is pruned: its bound $18$ cannot beat $20$. The optimum is $(4, 0)$ with value $20$, which enumerating all feasible integer points confirms; rounding the relaxed optimum to $(3, 2)$ would be infeasible and to $(3, 1)$ suboptimal.
\end{example}

\paragraph{Big-$M$ constraints.}
Binaries make it possible to say ``at least one of these constraints holds'',
which no LP can express. Given a constant $M$ larger than any value that
$\vect{a}\T\vect{z} - \beta$ can take on the feasible set, the pair
\begin{equation}
  \vect{a}\T\vect{z} \le \beta + M\,b,\qquad b \in \set{0, 1},
  \label{eq:appB-bigM}
\end{equation}
enforces the constraint when $b = 0$ and switches it off when $b =
1$.\index{big-M constraint} To keep a drone out of an axis-aligned box
$[x_{\min}, x_{\max}] \times [y_{\min}, y_{\max}]$, \cref{ch:ch22} writes the
four ``outside'' conditions $x \le x_{\min}$, $x \ge x_{\max}$, $y \le
y_{\min}$ and $y \ge y_{\max}$ in the form \cref{eq:appB-bigM} with binaries
$b_1, \dots, b_4$ and adds $b_1 + b_2 + b_3 + b_4 \le 3$, so that at least one
of the four conditions stays enforced and the position lies beyond at least
one face. The value of $M$ matters: too small cuts off feasible positions; too large
makes the LP relaxation so loose that branch and bound has little to prune,
besides the numerical trouble of mixing coefficients of very different sizes.

\subsection{Gradient descent}
\label{sec:appB-gd}

When a problem has no constraints and is too large, too nonlinear or too
irregular for the solvers above, the tool of last resort is to follow the
negative gradient:\index{gradient descent}
\begin{equation}
  \vect{x}_{k+1} = \vect{x}_k - \alpha\,\nabla f(\vect{x}_k),\qquad \alpha > 0 \text{ the step size (learning rate)}.
  \label{eq:appB-gd}
\end{equation}
By \cref{thm:appB-taylor}, $f(\vect{x}_{k+1}) \approx f(\vect{x}_k) -
\alpha\norm{\nabla f(\vect{x}_k)}^2$, so a small enough step always decreases
$f$; the question is how small.

\begin{proposition}[Convergence of gradient descent]
\label{thm:appB-gd}
Let $f$ be convex and differentiable with an $L$-Lipschitz gradient, $\norm{\nabla f(\vect{x}) - \nabla f(\vect{y})} \le L\norm{\vect{x} - \vect{y}}$, and let $\vect{x}^*$ be a minimiser. With a constant step $\alpha \le 1/L$ the iterates of \cref{eq:appB-gd} satisfy $f(\vect{x}_k) - f(\vect{x}^*) \le \norm{\vect{x}_0 - \vect{x}^*}^2/(2\alpha k)$. If in addition $f$ is $\mu$-strongly convex, $\nabla^2 f \succeq \mu\mat{I}$, then $\norm{\vect{x}_k - \vect{x}^*} \le (1 - \alpha\mu)^{k}\norm{\vect{x}_0 - \vect{x}^*}$ for $\alpha \le 1/L$: the error shrinks geometrically, at a rate governed by the condition number $L/\mu$.
\end{proposition}

Proofs are in \textcite[Chapter~9]{boyd2004convex} and
\textcite[Chapter~3]{nocedal2006numerical}. The quadratic $f(\vect{x}) =
\tfrac12\vect{x}\T\mat{A}\vect{x}$ with $\mat{A} \succ \mat{0}$ shows the
mechanism exactly: in the eigenbasis of $\mat{A}$ each coordinate is
multiplied by $1 - \alpha\lambda_i$ per step, so the iteration converges iff
$\abs{1 - \alpha\lambda_i} < 1$ for all $i$, that is iff $\alpha <
2/\lambda_{\max}$, and the slowest coordinate is the one with the smallest
eigenvalue. A large eigenvalue ratio, a long narrow valley, forces a small
step and a slow crawl along the valley floor.

\begin{example}[Step sizes on a narrow valley]
\label{ex:appB-gd}
Take $f(\vect{x}) = \tfrac12(x_1^2 + 10x_2^2)$, so $\lambda_{\min} = 1$, $\lambda_{\max} = 10$ and $\nabla f = (x_1, 10x_2)\T$, and start at $\vect{x}_0 = (9, 2)\T$. With $\alpha = 0.05$ the factors are $0.95$ and $0.5$: the iterates drop to the valley floor and then crawl along it, reaching $\vect{x}_{20} = (3.23, 0.00)\T$ with $f = 5.20$. With $\alpha = 0.18$ the factors are $0.82$ and $-0.80$: the iterates zigzag across the valley but make faster progress along it, $\vect{x}_{20} = (0.170, 0.023)\T$ with $f = 0.017$. With $\alpha = 0.21$ the factor $-1.1$ on $x_2$ makes the iteration diverge: after $20$ steps $x_2 = 13.5$ and $f = 905$. \Cref{fig:appB-gd} shows the first two paths.
\end{example}

\begin{figure}[tb]
  \centering
  \inputfigure{appB/gradient-descent}
  \caption{Gradient descent on $f = \tfrac12(x_1^2 + 10x_2^2)$ from $(9, 2)$ (\cref{ex:appB-gd}). The ellipses are level sets. The small step $\alpha = 0.05$ (blue) reaches the valley floor at once and then advances by only $5\,\%$ per step; the larger step $\alpha = 0.18$ (orange) zigzags across the valley and gets much closer to the minimum in the same $20$ steps. Any $\alpha > 0.2$ diverges.}
  \label{fig:appB-gd}
\end{figure}

The same picture explains two phenomena in the book. The potential-field
controller of \cref{ch:ch15} is gradient descent on its potential, with the
drone's motion as the iteration: its oscillations in narrow corridors are the
zigzag of \cref{fig:appB-gd}, its local minima the non-convexity of
\cref{fig:appB-convex}. Training the prediction networks of \cref{ch:ch20} is
gradient descent on a loss averaged over data, with the gradient obtained by
the chain rule \cref{eq:appB-chain} on random mini-batches (\emph{stochastic}
gradient descent); the learning rate is the $\alpha$ above, and momentum and
adaptive step sizes are remedies for the narrow-valley effect.

% ---------------------------------------------------------------------
\section{Graphs and the Laplacian}
\label{sec:appB-graphs}

\subsection{Adjacency, degree and Laplacian matrices}

An undirected graph $G = (V, E)$ with $n$ vertices numbered $1, \dots, n$ is
the communication network of \cref{ch:ch23}: the vertices are drones and an
edge $\set{i, j}$ means that $i$ and $j$ can exchange messages. The
neighbours of $i$ are $N_i = \set{j : \set{i, j} \in E}$.

\begin{definition}[Adjacency, degree and Laplacian matrices]
\label{def:appB-laplacian}
The \textbf{adjacency matrix}\index{adjacency matrix} $\mat{A} \in \R^{n \times n}$ has $a_{ij} = a_{ji} = 1$ if $\set{i, j} \in E$ and $0$ otherwise (for a weighted graph, $a_{ij} = w_{ij} \ge 0$). The \textbf{degree} of $i$ is $d_i = \sum_j a_{ij}$, the degree matrix is $\mat{D} = \diag(d_1, \dots, d_n)$, and the \textbf{graph Laplacian}\index{graph Laplacian} is
\begin{equation}
  \mat{L} = \mat{D} - \mat{A},\qquad
  (\mat{L}\vect{x})_i = \sum_{j \in N_i} a_{ij}(x_i - x_j).
  \label{eq:appB-laplacian}
\end{equation}
\end{definition}

The second form of \cref{eq:appB-laplacian} is the one to remember: applied to
a vector of drone states, $\mat{L}$ returns for every drone the sum of its
disagreements with its neighbours, which is exactly what a consensus
protocol feeds back.

\begin{proposition}[The Laplacian quadratic form]
\label{thm:appB-laplacian-form}
For every $\vect{x} \in \R^n$, $\displaystyle \vect{x}\T\mat{L}\vect{x} = \sum_{\set{i,j} \in E} a_{ij}\,(x_i - x_j)^2$.
\end{proposition}
\begin{proof}
$\vect{x}\T\mat{D}\vect{x} - \vect{x}\T\mat{A}\vect{x} = \sum_i d_i x_i^2 - \sum_{i,j} a_{ij}x_i x_j = \tfrac12\sum_{i,j} a_{ij}(x_i^2 + x_j^2 - 2x_i x_j) = \tfrac12\sum_{i,j} a_{ij}(x_i - x_j)^2$, using $d_i = \sum_j a_{ij}$; each edge appears twice in the double sum, which removes the factor $\tfrac12$. \qedhere
\end{proof}

\begin{theorem}[Spectrum of the Laplacian]
\label{thm:appB-laplacian}
Let $\mat{L}$ be the Laplacian of an undirected graph with eigenvalues $\lambda_1 \le \lambda_2 \le \dots \le \lambda_n$ and maximum degree $d_{\max}$.
\begin{enumerate}[label=(\roman*)]
  \item $\mat{L}$ is symmetric and PSD.
  \item $\mat{L}\vect{1} = \vect{0}$, so $\lambda_1 = 0$ with eigenvector $\vect{1}$.
  \item The multiplicity of the eigenvalue $0$ equals the number of connected components of $G$. In particular $\lambda_2 > 0$ iff $G$ is connected.
  \item $\lambda_n \le 2 d_{\max}$.
\end{enumerate}
\end{theorem}
\begin{proof}
(i) Symmetry is inherited from $\mat{A}$ and $\mat{D}$; $\vect{x}\T\mat{L}\vect{x} \ge 0$ by \cref{thm:appB-laplacian-form}. (ii) Every row of $\mat{L}$ sums to $d_i - \sum_j a_{ij} = 0$. (iii) By \cref{thm:appB-laplacian-form}, $\vect{x}\T\mat{L}\vect{x} = 0$ iff $x_i = x_j$ on every edge, that is, iff $\vect{x}$ is constant on each connected component; for a PSD matrix, $\vect{x}\T\mat{L}\vect{x} = 0$ iff $\mat{L}\vect{x} = \vect{0}$ (write $\vect{x}$ in the eigenbasis). The null space is therefore spanned by the indicator vectors of the components, which are linearly independent, so its dimension is the number of components, and this dimension is the multiplicity of $0$ for a symmetric matrix. (iv) By Gershgorin's theorem every eigenvalue lies in one of the discs centred at $l_{ii} = d_i$ with radius $\sum_{j \ne i}\abs{l_{ij}} = d_i$, hence in $[0, 2d_{\max}]$. \qedhere
\end{proof}

The eigenvalue $\lambda_2$ is the \textbf{algebraic
connectivity}\index{algebraic connectivity} of the graph
\cite{fiedler1973algebraic}: it is zero exactly when the network is cut, and
the larger it is, the faster information spreads. \Cref{ch:ch23} shows that
under the consensus dynamics $\dot{\vect{x}} = -\mat{L}\vect{x}$ the average
$\tfrac1n\vect{1}\T\vect{x}$ is conserved, because $\vect{1}\T\mat{L} =
\vect{0}\T$ by (i) and (ii), and the disagreement decays like $e^{-\lambda_2
t}$ \cite{olfatisaber2007consensus}. The discrete version $\vect{x}_{k+1} =
(\mat{I} - \eps\mat{L})\vect{x}_k$ multiplies the eigencomponent $i$ by $1 -
\eps\lambda_i$ per step, so it converges iff $0 < \eps < 2/\lambda_n$, and by
(iv) any $\eps < 1/d_{\max}$ is safe without knowing the spectrum.

\begin{example}[A four-drone network]
\label{ex:appB-laplacian}
The graph of \cref{fig:appB-laplacian} has the edges $\set{1,2}$, $\set{2,3}$, $\set{3,4}$ and $\set{2,4}$, so the degrees are $(1, 3, 2, 2)$ and
\begin{equation*}
  \mat{L} = \begin{bmatrix} 1 & -1 & 0 & 0\\ -1 & 3 & -1 & -1\\ 0 & -1 & 2 & -1\\ 0 & -1 & -1 & 2\end{bmatrix},
  \qquad \lambda = (0, 1, 3, 4).
\end{equation*}
The graph is connected, so $\lambda_2 = 1 > 0$, and $\lambda_n = 4 \le 2d_{\max} = 6$. The eigenvector of $\lambda_2$, the Fiedler vector $(-2, 0, 1, 1)\T$, is negative on drone $1$ and positive on drones $3$ and $4$: its sign pattern points at the bridge $\set{1,2}$ as the weakest link. Deleting that edge isolates drone $1$ and the spectrum becomes $(0, 0, 3, 3)$, with the double zero of \cref{thm:appB-laplacian}(iii). From $\vect{x}_0 = (4, 0, 1, -1)\T$ the consensus dynamics converge to the average $1$ on every drone.
\end{example}

\begin{figure}[tb]
  \centering
  \inputfigure{appB/laplacian}
  \caption{The network of \cref{ex:appB-laplacian} and its Laplacian. Each diagonal entry is a degree and each $-1$ an edge. The entries of the Fiedler vector are written next to the vertices; the sign change across the edge $\set{1,2}$ marks the cut that disconnects the graph.}
  \label{fig:appB-laplacian}
\end{figure}

\subsection{Shortest-path notation}
\label{sec:appB-shortest-path}

The search chapters use a directed graph $G = (V, E)$ with a cost $c(u, v)
\ge 0$ on every edge, successors $\Succ(u) = \set{v : (u, v) \in E}$ and
predecessors $\Pred(u)$. A path $\pi = (v_0, \dots, v_m)$ follows edges and
costs $\cost(\pi) = \sum_i c(v_i, v_{i+1})$; the \textbf{shortest-path
distance}\index{shortest-path distance} $\dist(u, v)$ is the least cost of a
path from $u$ to $v$, $\infty$ if there is none, and it satisfies the triangle
inequality $\dist(u, w) \le \dist(u, v) + \dist(v, w)$. \Cref{ch:ch03}
computes $\dist(s, \cdot)$ from a start $s$ and writes it $\delta(s, v)$;
\cref{ch:ch04} keeps for every node the cost $\gcost(n)$ of the best path
found so far, a heuristic estimate $\hcost(n)$ of $\dist(n, \gamma)$ to the
goal $\gamma$ and the priority $\fcost = \gcost + \hcost$, with the frontier in
\Open and the expanded nodes in \Closed. A 4- or 8-connected grid is such a
graph with $\abs{E} \le 4\abs{V}$ or $8\abs{V}$, and the space-time graph of
\cref{ch:ch04,ch:ch08} has the $(T+1)\abs{V}$ states $(v, t)$, $t = 0, \dots,
T$, with a wait edge from $(v, t)$ to $(v, t+1)$ beside the move edges. 

% ---------------------------------------------------------------------
\section{Complexity and data structures}
\label{sec:appB-complexity}

\subsection{Asymptotic notation}

\begin{definition}[Big-O, Omega and Theta]
\label{def:appB-bigO}
For functions $f, g\colon \N \to \R_{\ge 0}$, $f(n) = \bigO{g(n)}$ if there are constants $c > 0$ and $n_0$ with $f(n) \le c\,g(n)$ for all $n \ge n_0$; $f(n) = \Omega(g(n))$ if $f(n) \ge c\,g(n)$ for all $n \ge n_0$; and $f(n) = \Theta(g(n))$ if both hold.\index{big-O notation}
\end{definition}

Big-O is an upper bound up to a constant factor for large inputs; $\Theta$
says the bound is tight. Constants and lower-order terms are dropped ($3n^2 +
50n = \Theta(n^2)$), the base of a logarithm is irrelevant, and the usual
ladder is $1 \prec \log n \prec n \prec n\log n \prec n^2 \prec 2^n \prec n!$.
A bound describes the worst case unless it says \emph{expected} (average over
random choices) or \emph{amortised} (average over a sequence of operations).
Read a bound with the right variables: Dijkstra's algorithm with a binary heap
runs in $\bigO{(\abs{V} + \abs{E})\log\abs{V}}$ (\cref{ch:ch03}), which on a
grid with $\abs{E} = \Theta(\abs{V})$ is $\bigO{\abs{V}\log\abs{V}}$; a
space-time search over $T$ steps has $(T+1)\abs{V}$ states; and the
exponential worst case of the multi-agent solvers of \cref{ch:ch09,ch:ch10} is
in the number of agents, not in the size of the map.

\subsection{Binary heaps}

A \textbf{binary heap}\index{binary heap} stores $n$ keys in an array $a[0],
\dots, a[n-1]$ read as a complete binary tree: the children of position $i$
are $2i+1$ and $2i+2$, its parent is $\lfloor (i-1)/2 \rfloor$, and the
\emph{heap property} says that no key is larger than its children's, so
$a[0]$ is the minimum. Inserting appends the key and swaps it upwards while it
is smaller than its parent (\emph{sift-up}); extracting the minimum moves the
last key to the root and swaps it downwards with its smaller child
(\emph{sift-down}); each walks one root-to-leaf path of length $\lfloor \log_2
n \rfloor$, hence $\bigO{\log n}$. Building a heap from an unsorted array by
sifting down from the last internal node upwards costs $\bigO{n}$
\cite[Chapter~6]{cormen2009clrs}. \Cref{tab:appB-heap} collects the costs of
the operations a search needs, together with those of the \python containers
used in this book's code.

\begin{table}[tb]
  \centering\small
  \caption{Costs of the operations of a binary heap and of the \python containers used in the search code. Hash-table costs are expected costs; worst cases are $\bigO{n}$.}
  \label{tab:appB-heap}
  \begin{tabular}{@{}llll@{}}
  \toprule
  Structure & Operation & Cost & \python\\
  \midrule
  binary heap & peek minimum & $\bigO{1}$ & \code{heap[0]}\\
  binary heap & insert & $\bigO{\log n}$ & \code{heapq.heappush}\\
  binary heap & extract minimum & $\bigO{\log n}$ & \code{heapq.heappop}\\
  binary heap & build from $n$ items & $\bigO{n}$ & \code{heapq.heapify}\\
  binary heap & decrease-key & $\bigO{\log n}$ & push again; skip stale entries\\
  hash table & insert, lookup, delete & $\bigO{1}$ exp. & \code{dict}, \code{set}\\
  dynamic array & append; index & $\bigO{1}$ amort.; $\bigO{1}$ & \code{list.append}, \code{a[i]}\\
  dynamic array & pop front; membership & $\bigO{n}$ & \code{list.pop(0)}, \code{x in a}\\
  sorted array & insert & $\bigO{n}$ & \code{bisect.insort}\\
  any & sort $n$ items & $\bigO{n\log n}$ & \code{sorted}\\
  \bottomrule
  \end{tabular}
\end{table}

\subsection{Hashing, and the \code{heapq} and \code{dict} idiom}

A \textbf{hash table}\index{hash table} maps a key to an array position with
a hash function and resolves the occasional collision by probing or chaining,
so that insertion, lookup and deletion take constant expected time
\cite[Chapter~11]{cormen2009clrs}; \code{dict} and \code{set} are hash tables.
Keys must be \emph{hashable}, which in \python means immutable: tuples of
integers such as a cell $(x, y)$ or a space-time state $(x, y, t)$ are ideal,
lists and NumPy arrays are not accepted, and floating-point keys are legal but
dangerous because two values that are equal on paper may differ in the last
bit. This is why the search code of the book indexes cells by integer tuples.

The \code{heapq} module implements a binary heap on a plain list, but it has
no decrease-key. The idiom of \cref{ch:ch03,ch:ch04}, inherited by every
planner built on them, is \emph{lazy deletion}\index{lazy deletion}: keep the
best known $\gcost$ of every node in a \code{dict}; when a better cost is
found, push a fresh entry $(\fcost, -\gcost, \mathit{counter}, \mathit{node})$
and leave the old one in the heap; when an entry is popped, skip it if its node
is already in the closed set. The heap then holds at most one entry per
relaxation, $\bigO{\abs{E}}$ in total, which leaves the asymptotic bounds
unchanged, and the \code{dict} lookup that guards each push costs constant
expected time.

\begin{pitfall}[Heap entries must be comparable]
\code{heapq} orders tuples lexicographically, so when two entries agree on $\fcost$ and $\gcost$ it compares the next field. If that field is a node without an ordering (a dictionary, an object) the search crashes at the first tie, and if it is a NumPy array it crashes at the first comparison. A running counter from \code{itertools.count()}, placed before the node, guarantees that the comparison never reaches the node and makes tie-breaking deterministic; the entry $-\gcost$ before the counter prefers the larger $\gcost$ among equal $\fcost$, as \cref{ch:ch04} recommends.
\end{pitfall}

% ---------------------------------------------------------------------
\section{Summary}
\begin{summary}
\begin{itemize}
  \item A symmetric matrix is a scaling along orthogonal axes, $\mat{A} = \mat{Q}\mat{\Lambda}\mat{Q}\T$; it is PSD iff all $\lambda_i \ge 0$, and every covariance is PSD. Block matrices are handled with the Schur complement, which gives the block inverse, the block determinant and the Woodbury identity.
  \item Least squares is solved by the normal equations $\mat{A}\T\mat{A}\vect{x} = \mat{A}\T\vect{b}$; never form an inverse to do it.
  \item Linearisation keeps the Jacobian and drops an error of second order in the displacement; the double-integrator step $\pos' = \pos + \vel\dt + \tfrac12\acc\dt^2$, $\vel' = \vel + \acc\dt$ is exact for a piecewise-constant input.
  \item Covariances propagate as $\mat{A}\mat{\Sigma}\mat{A}\T$ under any linear map; Gaussians stay Gaussian under affine maps, marginalisation and conditioning, and the conditioning formula \cref{eq:appB-conditional} is the Kalman update in disguise.
  \item A filter is total probability (predict) followed by Bayes' rule (update); Monte Carlo errors shrink like $1/\sqrt{N}$; the squared Mahalanobis distance of a Gaussian vector is $\chi^2_n$, which sets gates such as $5.991$ for $n = 2$ at $95\,\%$.
  \item Convexity makes local minima global and the KKT conditions sufficient. An LP attains its optimum at a vertex, a convex QP is an LP with a quadratic cost, a MILP is solved exactly by branch and bound with big-$M$ binaries encoding disjunctions, and gradient descent converges for $\alpha < 2/\lambda_{\max}$ but crawls on ill-conditioned problems.
  \item The Laplacian $\mat{L} = \mat{D} - \mat{A}$ is PSD with $\mat{L}\vect{1} = \vect{0}$; $\lambda_2 > 0$ iff the graph is connected and sets the speed of consensus, and $\lambda_n \le 2d_{\max}$ bounds the safe step size.
  \item A binary heap gives $\bigO{\log n}$ insertion and extraction; \code{heapq} plus a \code{dict}, lazy deletion and a tie-breaking counter is the priority queue of every search in this book.
\end{itemize}
\end{summary}

\paragraph{Further reading.}
Linear algebra, convexity and optimality conditions are treated with great
clarity in \textcite{boyd2004convex}, whose appendices cover the matrix facts
of \cref{sec:appB-linalg}; the algorithms of optimization are in
\textcite{nocedal2006numerical} and integer programming in
\textcite{wolsey1998integer}. The probabilistic material follows
\textcite{thrun2005probabilistic} and \textcite{sarkka2013bayesian}; gating is
treated in \textcite{barshalom2001estimation} and sequential Monte Carlo in
\textcite{doucet2001sequential}. Spectral graph theory and consensus are
surveyed in \textcite{olfatisaber2007consensus}, building on
\textcite{fiedler1973algebraic}; heaps, hashing and asymptotic notation are in
\textcite{cormen2009clrs}.

% ---------------------------------------------------------------------
\section{Exercises}
\label{sec:appB-exercises}

\begin{exercise}[\difficulty{1}]
\label{exr:appB-psd}
Show that the sum of two PSD matrices is PSD, and decide whether $\begin{bsmallmatrix} 1 & 2\\ 2 & 1\end{bsmallmatrix}$ can be a covariance matrix.
\end{exercise}

\begin{exercise}[\difficulty{2}]
\label{exr:appB-woodbury}
Apply \cref{thm:appB-woodbury} to $\mat{P} - \mat{P}\mat{H}\T(\mat{H}\mat{P}\mat{H}\T + \mat{R})\inv\mat{H}\mat{P}$ and show that it equals $(\mat{P}\inv + \mat{H}\T\mat{R}\inv\mat{H})\inv$. Which form is cheaper when there are many more state variables than measurement components?
\end{exercise}

\begin{exercise}[\difficulty{2}]
\label{exr:appB-conditioning}
In \cref{ex:appB-conditioning}, condition on $x = 3$ instead and find the distribution of $y$. Check that the variance shrinks by the same factor $1 - \rho^2$.
\end{exercise}

\begin{exercise}[\difficulty{3} Coding]
\label{exr:appB-coding}
Starting from \code{code/figures/gen\_appB\_numbers.py}, draw $N = 10\,000$ samples from the Gaussian of \cref{ex:appB-conditioning}, estimate the mean and the covariance from the samples, and count the fraction of samples whose squared Mahalanobis distance exceeds the $95\,\%$ gate $5.991$ for $n = 2$. Then perturb the covariance used in the gate by $\pm 30\,\%$ and report how the rejection rate changes.
\end{exercise}

\index{mathematical refresher|)}
